% Revised October 29, 2026 from the supplied psh_exhaustion.tex.
% The source was compiled and checked against the supplied 28-page PDF.
\documentclass[11pt]{amsart}

\usepackage[T1]{fontenc}
\usepackage{lmodern}
\usepackage{amssymb,mathtools}
\usepackage{microtype}
\usepackage{enumitem}
\usepackage{booktabs,tabularx,placeins}
\usepackage{xcolor}
\usepackage[colorlinks=true,linkcolor=blue!55!black,citecolor=blue!55!black,urlcolor=blue!55!black]{hyperref}

\newtheorem{theorem}{Theorem}[section]
\newtheorem{proposition}[theorem]{Proposition}
\newtheorem{lemma}[theorem]{Lemma}
\newtheorem{corollary}[theorem]{Corollary}
\theoremstyle{definition}
\newtheorem{definition}[theorem]{Definition}
\theoremstyle{remark}
\newtheorem{remark}[theorem]{Remark}
\numberwithin{equation}{section}

\newcommand{\C}{\mathbb C}
\newcommand{\R}{\mathbb R}
\newcommand{\PP}{\mathbf P}
\newcommand{\OO}{\mathcal O}
\newcommand{\Tot}{\operatorname{Tot}}
\newcommand{\Bl}{\operatorname{Bl}}
\newcommand{\ddc}{dd^c}
\newcommand{\psh}{\mathrm{psh}}
\newcommand{\Levi}{\mathcal L}
\newcommand{\barp}{\bar\partial}
\newcommand{\supp}{\operatorname{supp}}
\newcommand{\base}{\operatorname{pr}}
\newcommand{\Uone}{\mathrm U(1)}

\title[Continuous and smooth plurisubharmonic exhaustions]{Continuous and Smooth
Plurisubharmonic Exhaustions:\\
Threefold Counterexamples and Surface Regularization}
\author{Zhijian Chen}
\address{Institute of Mathematics, Academy of Mathematics and Systems Science, Chinese Academy of Sciences, Beijing 100190, China}
\email{chenzhijian@amss.ac.cn}
\author{Bo Xiao}
\address{School of Mathematical Sciences, University of Chinese Academy of Sciences, Beijing, China}
\email{xb1994@amss.ac.cn}
\date{October 29, 2026}
\subjclass[2020]{Primary 32F17; Secondary 32Q15, 32L05, 32U40}
\keywords{Plurisubharmonic exhaustion, K\"ahler manifold, semipositive line bundle,
ruled surface, Hironaka modification, Ueda theory, class VII surface, boundary divisor, positive current}

\begin{document}


\begin{abstract}
Starting from Koike's modified Grauert example, we construct a smooth
quasi-projective K\"ahler threefold with a locally Lipschitz
plurisubharmonic exhaustion but no $C^2$ plurisubharmonic exhaustion.
A modification over one point gives a non-K\"ahler threefold with the
same properties.  In dimension two, we prove that every divisor
complement in a compact K\"ahler surface admitting a continuous
plurisubharmonic exhaustion admits a smooth one.  In particular, this
holds for every smooth complex quasi-projective surface.  The proof
averages positive $dd^c$-closed currents over exhaustion levels to
eliminate their Lelong numbers, then combines the surface current
theorem of Dinh--Nguy\^en--Sibony with the Hodge index theorem and
the intersection matrix of the boundary.  In the non-Stein-modification
case, a positive integral divisor supported on the entire boundary has
a flat Hermitian metric near the boundary and yields a pluriharmonic
logarithmic collar.  We retain numerical existence criteria for
nonprojective compactifiers and smooth rulings, and sufficient
conditions from Ueda theory.  The compactification hypotheses and the
recent elliptic-neighborhood input to these separate existence results
are explicit.  In arbitrary dimension, we also show that a reduced
pure $n$-dimensional complex space with a continuous psh exhaustion
and a holomorphic map to $\C^{n-1}$ whose fibers all have dimension
at most one is holomorphically convex, hence admits a smooth psh
exhaustion, without any restriction on its singular locus.  This
criterion combines geometric flatness and cycle compactness with
V\^{\i}j\^{\i}itu's convexity criteria.  The unrestricted regularization
problem for noncompact complex surfaces is not resolved here.
\end{abstract}

\maketitle

\section{Introduction and main results}\label{sec:intro}

The regularity of a plurisubharmonic exhaustion matters even for
K\"ahler manifolds.  This paper constructs a smooth quasi-projective
threefold admitting a continuous plurisubharmonic exhaustion but no
$C^2$ one, and then inserts a non-K\"ahler configuration over a single
point without changing these exhaustion properties.  In contrast, we
prove continuous-to-smooth regularization for divisor complements in
compact K\"ahler surfaces, including every smooth quasi-projective
surface.  A further affirmative result in every dimension uses
holomorphic maps with one-dimensional fibers and also permits a
singular total space, with no restriction on its singular locus.
Throughout the introduction, psh abbreviates plurisubharmonic,
and smooth means $C^\infty$.

\subsection{Threefold counterexamples and their metric construction}

Our compact base is the modified Grauert ruled surface appearing in
Koike's work \cite[Example~4.1]{Koike2021}. Its failure of semipositivity
is already known; the first-order obstruction, including its $C^2$
version, is not claimed as a new result here
\cite[Theorem~1.4 and Example~4.1]{Koike2021},
\cite[Remark~2.8]{Wu2026}. We include a direct proof to make the
regularity threshold used below explicit. The analytic construction
then consists of a cusp correction, a global maximum, and a radius
criterion converting any hypothetical $C^2$ exhaustion into a
forbidden metric. Failure of smooth semipositivity by itself would
not provide a continuous exhaustion.

The cusp estimate also reaches the locally Lipschitz endpoint: a
sufficiently large multiple of $q^{1/2}$ dominates the initial
curvature in the scaled tangent variables. The coefficient is
essential at this endpoint; for smaller exponents shrinking the
neighborhood already suffices.

\begin{theorem}\label{thm:main}
Let $C$ be the smooth projective completion of
\[
 y^2=x^5-x,
\]
let $K=K_C$, and let
\begin{equation}\label{eq:extension-intro}
 0\longrightarrow K\longrightarrow V\longrightarrow\OO_C
 \longrightarrow0
\end{equation}
be a nonsplit holomorphic extension. Let
\[
 S=\PP_{\mathrm{lines}}(V),\qquad
 \pi:S\longrightarrow C,\qquad
 Y=\PP(K)\subset S,
\]
and set
\begin{equation}\label{eq:main-line}
 L=\OO_S(Y)\otimes\pi^*K,\qquad M=\Tot(L^{-1}).
\end{equation}
Then:
\begin{enumerate}[label=\textup{(\roman*)},leftmargin=2.5em]
 \item $S$ is a smooth projective surface and $M$ is a smooth
 quasi-projective K\"ahler threefold;
 \item $L$ admits a continuous Hermitian metric with plurisubharmonic
 local weights, but admits no $C^2$ Hermitian metric with semipositive
 curvature;
 \item $M$ admits a continuous plurisubharmonic exhaustion, but admits
 no $C^2$ plurisubharmonic exhaustion.
\end{enumerate}
The metric and the exhaustion in \textup{(ii)--(iii)} can be chosen
locally Lipschitz, and hence locally of class $C^{0,\gamma}$ for every
$0<\gamma\le1$. No assertion about existence of a $C^1$ exhaustion
is made.
\end{theorem}

The passage from metrics to exhaustions uses a criterion valid for a
holomorphic line bundle $A$ over any compact complex manifold $Z$:
\[
\begin{gathered}
 \Tot(A^{-1})\text{ admits a }C^r\text{ psh exhaustion}\\
 \Longleftrightarrow\\[-2pt]
 A\text{ admits a }C^r\text{ semipositive Hermitian metric},
 \qquad r\in\{2,\infty\}.
\end{gathered}
\]
This is Corollary~\ref{cor:criterion}. A semipositive metric gives the
squared dual norm exhaustion. Conversely, averaging an arbitrary
exhaustion under fiberwise circle rotation and taking a sufficiently
high level yields a radius function whose negative logarithm is a psh
metric weight. Consequently, the obstruction in Theorem~\ref{thm:main}
rules out every $C^2$ psh exhaustion, not only a smoothing of the
particular one constructed here.

The non-K\"ahler example has the same dimension and agrees with $M$
away from a compact exceptional fiber.

\begin{theorem}\label{thm:nonkahler-threefold}
For every point $x$ of the threefold $M$ in Theorem~\ref{thm:main}, there
is a smooth non-K\"ahler threefold $N$ and a proper holomorphic
modification $\mu:N\to M$ inducing a biholomorphism
\[
 N\setminus\mu^{-1}(x)\simeq M\setminus\{x\}.
\]
The threefold $N$ admits a locally Lipschitz psh exhaustion but no
$C^2$ psh exhaustion.
\end{theorem}

The construction places Hironaka's opposite-order blowups inside the
exceptional plane of the blowup of $x$. It produces two compact rational
curves whose sum is zero in real homology, precluding a K\"ahler form.
The exhaustion assertions follow from Proposition~\ref{prop:compact-modification}:
a proper surjective holomorphic map that is biholomorphic away from a compact
subset preserves the existence of $C^r$ psh exhaustions, for
$r=0,2,\infty$. The reverse implication uses convex composition to
make an exhaustion constant near the exceptional set before descending
it. Thus the modification is localized over one point; no small
deformation of a fixed complex structure is asserted. Choosing $x$
outside the zero section also preserves a neighborhood of the compact
base. The compact product principle in Proposition~\ref{prop:compact-product}
gives further examples in higher dimensions, including the Hopf-surface
product in Corollary~\ref{cor:nonkahler}.


\subsection{Regularization on surfaces with a divisor compactification}
The surface result has a different conclusion from the threefold
construction.

\begin{theorem}\label{thm:surface-introduction}
Let $S$ be a smooth compact connected K\"ahler surface and let
$D\ne0$ be an effective divisor.  If
$X=S\setminus\supp D$ admits a continuous psh exhaustion, then
it admits a smooth psh exhaustion.  In particular, every smooth
complex quasi-projective surface admitting a continuous psh
exhaustion admits a smooth one.
\end{theorem}

Theorem~\ref{thm:divisor-regularization} and
Corollary~\ref{cor:quasiprojective-regularization} prove this statement
without assuming that $S$ is minimal or that the boundary is connected,
smooth, or has normal crossings.  The proof first uses the continuous
minimal-kernel criterion of Mongodi--S{\l}odkowski
\cite[Lemma~2.1]{MongodiSlodkowski2020}.  Unless $X$ is a modification of a Stein
space, sufficiently high exhaustion levels support positive
$dd^c$-closed currents.  Finite averages over equally spaced levels
have a mass-one limit supported compactly inside $X$.  A uniform
Lelong--Jensen estimate and continuity of the exhaustion force all
Lelong numbers of this limit to vanish.  The surface current theorem
of Dinh--Nguy\^en--Sibony \cite[Corollary~2.4 and its proof]{DNS22}
then combines with the Hodge index theorem and the boundary
intersection matrix to produce a positive integral divisor $A$ with
full boundary support, $A^2=0$, and a smooth flat Hermitian metric on
$\mathcal O_S(A)$ near the boundary.  Its canonical section yields
a smooth pluriharmonic logarithmic collar.

In particular, a continuous divergent psh collar around a smooth
square-zero curve in a compact K\"ahler surface forces a smooth
exhaustion, with no restriction on the genus or Ueda type
(Corollary~\ref{cor:square-zero-continuous-collar}).  The proof
does not infer convergence merely from infinite Ueda type.  Flatness
is obtained in the non-Stein-modification alternative by the current
and intersection argument.

Section~5 also retains results that establish existence without
assuming a continuous exhaustion: numerical boundary criteria for
nonprojective compactifiers and smooth rulings, and sufficient
conditions from Ueda theory.  These are logically distinct from
regularization.  For arbitrary noncompact surfaces with a nonconstant
holomorphic function, V\^{\i}j\^{\i}itu's holomorphic-convexity theorem
\cite{Vajaitu2022} gives the familiar regularization consequence
recorded in Proposition~\ref{prop:surface-nonconstant-function}.
Quasi-projectivity does not ensure the existence of such a function,
even under the exhaustion hypothesis; see
Remark~\ref{rem:quasiprojective-constant-functions}.  For a
line bundle over a compact curve, the degree criterion in
Proposition~\ref{cor:curve-total-space} gives an elementary
explanation of why the threefold construction cannot be repeated
over a curve.

Table~\ref{tab:surface-status} records the resulting
classification of the \emph{compactifiers} under a continuous
exhaustion hypothesis.  Every projective row is affirmative, including
nonminimal projective compactifiers.  An arbitrary noncompact
K\"ahler surface need not admit a divisor compactification, and an
arbitrary analytic domain in a projective surface need not have a
divisorial boundary.  Neither class is covered by the compactification
hypothesis.  Sections~2--4 contain the threefold constructions;
Section~\ref{sec:surfaces} establishes the independent surface results;
Section~\ref{sec:curve-fibers} proves the curve-fibering criterion
described next.

\subsection{Holomorphic convexity from curve fiberings}
An additional regularization mechanism requires no compactification
and works in every dimension.

\begin{theorem}\label{thm:curve-fiber-introduction}
Let $X$ be a reduced complex space of pure dimension $n\ge2$ with
a continuous real-valued psh exhaustion.  If a holomorphic map
\[
 F:X\longrightarrow\C^{n-1}
\]
has every nonempty fiber of dimension at most one, then $X$ is
holomorphically convex and admits a smooth psh exhaustion.
No condition is imposed on the connected components of
$\operatorname{Sing}X$.
\end{theorem}

Theorem~\ref{thm:curve-fiber-convexity} proves this assertion.
The dimension bound concerns every fiber, including exceptional
fibers; a generic rank condition would not suffice for the argument.
Normality or Cohen--Macaulayness of $X$, properness of $F$, and
smoothness or reduced scheme structure of its fibers are not assumed.
The smooth exhaustion is obtained by pulling back one from the Stein
target of the Remmert reduction.  Smoothness on a complex space means
local restriction of smooth ambient functions; the conclusion is
existence of a smooth exhaustion, not approximation of the given one.

The key step is the stability of noncompact irreducible fiber
components meeting a fixed compact set
(Lemma~\ref{curve:lem:stability}).
Geometric flatness over the normal base supplies weighted fiber
cycles even for a singular source.  The psh exhaustion confines
compact components to a common compact set, and Bishop compactness
then controls their limits.  This replaces the smooth-source
stability result in \cite[Corollary~4.2]{VijituForum2022};
the boundary-collar and sublevel criteria in that paper already
apply to complex spaces.  Corollary~\ref{cor:curve-contraction}
identifies the resulting simultaneous contraction of compact
fiber chains.

\subsection{Relation to prior work and scope of the claims}\label{subsec:provenance}
The known ruled-surface obstruction, Hironaka's opposite-order blowups,
and the quoted classification and Ueda theorems are inputs, not new
constructions of this paper.  The self-contained threefold argument
combines the explicit bounded-potential metric, the radius obstruction,
and a compactly supported modification.  The bounded-potential remark
at the end of Section~2 explains which part of the question in
\cite[Section~4]{Zhang2025} this construction addresses.

The compactifiable K\"ahler surface theorem combines the published
continuous minimal-kernel criterion and the general positive-current
theorem with the level-averaging estimate and the full-boundary
intersection argument given in Section~5.  It uses no elliptic
neighborhood theorem, surface abundance theorem, or minimal-model
reduction.  Within smooth quasi-projective manifolds, the examples
in Theorem~\ref{thm:main} occur in the smallest possible
dimension: surfaces are covered by
Theorem~\ref{thm:surface-introduction}, while a connected
noncompact Riemann surface is Stein.  We do not assert minimality
of dimension three among arbitrary noncompact K\"ahler manifolds,
or optimality of Lipschitz regularity as a threshold.

The elliptic-neighborhood step in the separate existence results
invokes \cite[Theorem~1.6]{CFSTZ2026}, a September~2026 preprint.
That input is independent of both threefold counterexamples and of
Theorem~\ref{thm:surface-introduction}.  No classification or general
regularization theorem for arbitrary noncompact complex surfaces
is asserted.

The curve-fibering criterion is a deduction from classical geometric
flatness and cycle compactness together with
\cite[Corollary~2.5 and Propositions~2.2, 2.4]{VijituForum2022}.
Its singular-source formulation does not require the compactness
condition on the singular locus in \cite[Theorem~1.1]{VijituForum2022}.
For an irreducible surface, its convexity conclusion already follows
from \cite[Theorem~1.1]{Vajaitu2022}.  We attribute the underlying
convexity and cycle-theoretic inputs explicitly and do not claim
priority for this formulation.

\subsection{Conventions and terminology}\label{sec:conventions}

We use
\[
 d^c=\frac{i}{4\pi}(\barp-\partial),\qquad
 \ddc=\frac{i}{2\pi}\partial\barp.
\]
For a holomorphic line bundle $A$ and a local holomorphic frame $e$, a
Hermitian metric is written
\[
 \|e\|_h^2=e^{-\varphi}.
\]
Its curvature is $\Theta_h(A)=\ddc\varphi$, representing $c_1(A)$.
If $e_j=g_{ij}e_i$, the metric weights satisfy
\begin{equation}\label{eq:metric-transition}
 \varphi_j=\varphi_i-\log|g_{ij}|^2.
\end{equation}
A positive continuous Hermitian metric is semipositive in the current
sense when its local weights are continuous psh functions. For a $C^2$
metric this agrees with pointwise semipositivity of its curvature.

For a $C^2$ real-valued function $f$, its unnormalized Levi form is
\[
 \Levi_f(\xi)=\sum_{j,k}f_{z_j\bar z_k}\xi_j\overline{\xi_k}.
\]
Thus $f$ is psh exactly when $\Levi_f\ge0$. The fixed positive factor
$1/(2\pi)$ in $\ddc$ does not affect any positivity argument.
Poincar\'e--Lelong takes the form
\begin{equation}\label{eq:PL}
 \ddc\log\|s_D\|_h^2=[D]-\Theta_h(\OO(D)),
\end{equation}
where $s_D$ is the canonical section of an effective divisor $D$.

\begin{definition}
A continuous real-valued function $F$ on a complex manifold $X$ is an
\emph{exhaustion} if every sublevel set $\{F\le c\}$ is compact.
Equivalently, $F(x_j)\to+\infty$ for every sequence escaping every
compact subset of $X$.
\end{definition}

We say that $X$ is \emph{continuously weakly complete} or
\emph{smoothly weakly complete} when it admits a continuous or smooth
psh exhaustion, respectively. Specifying the regularity avoids the
varying conventions for weak completeness in the literature.
In Section~\ref{sec:curve-fibers} the same terminology is used for
reduced complex spaces, with smooth functions understood through
local embeddings.
Our projectivization $\PP_{\mathrm{lines}}(E)$ parametrizes
one-dimensional subspaces of the fibers of $E$, and its tautological
line bundle is $\OO_{\PP(E)}(-1)$.

\section{The ruled surface and continuous semipositive metrics}\label{sec:metrics}

\subsection{The ruled surface and the metric obstruction}\label{sec:ruled}

\subsubsection{The curve and the extension}

The polynomial $x^5-x$ has five distinct roots. The double covering
$C\to\PP^1$ is therefore branched at these five roots and at infinity.
Riemann--Hurwitz gives
\[
 2g(C)-2=2(-2)+6=2,
\]
so $g(C)=2$ and $\deg K=2$. Serre duality gives
\[
 \operatorname{Ext}^1(\OO_C,K)=H^1(C,K)
 \simeq H^0(C,\OO_C)^*\simeq\C.
\]
Consequently a nonzero extension class exists. All nonzero classes give
isomorphic extensions after allowing scalar automorphisms of the end
terms, so the isomorphism class of $V$ is unambiguous.

One can also specify the extension through a Dolbeault operator. Choose
a smooth positive $(1,1)$-form $\beta$ on $C$ with
$\int_C\beta=1$, and regard it as a $K$-valued $(0,1)$-form. On the
smooth bundle $K\oplus\OO_C$, set
\begin{equation}\label{eq:dolbeault-extension}
 \barp_V(s,f)=(\barp_Ks+f\beta,\barp f).
\end{equation}
This operator is integrable, since there are no $(0,2)$-forms on a
curve, and defines \eqref{eq:extension-intro}. A splitting would give a
holomorphic lift $(s,1)$ and hence $\barp_Ks=-\beta$. Since $s$ is a
smooth $(1,0)$-form, Stokes' theorem would then imply
$\int_C\beta=0$, a contradiction. The choice of the standard
identification of $K$-valued $(0,1)$-forms with $(1,1)$-forms only changes
this argument by a sign.

\subsubsection{The negative section and the line bundle}

By GAGA, $V$ is algebraic, so $S=\PP_{\mathrm{lines}}(V)$ is a smooth
projective surface. The inclusion $K\hookrightarrow V$ gives a section
$Y=\PP(K)$ of $\pi$. Its normal bundle is
\begin{equation}\label{eq:normal}
 N_{Y/S}\simeq\operatorname{Hom}(K,V/K)\simeq K^{-1}.
\end{equation}
In particular,
\begin{equation}\label{eq:restriction}
 \OO_S(Y)|_Y\simeq K^{-1},\qquad
 L|_Y\simeq\OO_Y,\qquad Y^2=-2.
\end{equation}
The trivialization of $L|_Y$ in \eqref{eq:restriction} is the canonical
one induced by $K^{-1}\otimes K$.

For context, this line bundle is nef and big. If $F$ is a fiber of $\pi$,
then
\[
 L\cdot Y=0,\qquad L\cdot F=1,\qquad L^2=2.
\]
For every irreducible curve $D\ne Y$ dominating $C$,
\[
 L\cdot D=Y\cdot D+2\deg(\pi|_D)>0.
\]
Every vertical irreducible curve is a fiber. Hence $L$ is nef, and the
usual surface criterion $L^2>0$ implies bigness. For completeness, the
latter also follows from Riemann--Roch: $h^2(S,L^m)=0$ for $m$ large,
because $(K_S-mL)$ has negative intersection with a fixed ample divisor,
and therefore
\[
 h^0(S,L^m)\ge\chi(S,L^m)=m^2+O(m).
\]
Moreover, $2L-Y=Y+2\pi^*K$ is ample by Nakai--Moishezon: its square is
$6$, its intersections with $Y$ and a fiber are $2$ and $1$, and its
intersection with every other irreducible curve is positive.

\subsubsection{First-order transition functions}

Choose a finite coordinate cover $\{U_i\}$ of $C$, frames $k_i$ of $K$,
and holomorphic lifts $v_i$ of $1$ in \eqref{eq:extension-intro}.
On overlaps,
\begin{equation}\label{eq:vector-transitions}
 k_j=g_{ij}k_i,\qquad v_j=v_i+c_{ij}k_i.
\end{equation}
The cocycle $\{c_{ij}k_i\}$ represents the nonzero extension class.
Near $Y$, write a line in $V$ as
\[
 [k_i+w_iv_i].
\]
Expressing this line in the $j$-th frame yields
\begin{equation}\label{eq:normal-transition}
 w_j=\frac{g_{ij}w_i}{1-c_{ij}w_i}.
\end{equation}
Shrinking the neighborhoods if necessary, all denominators are nonzero.
Let $d_i$ be the frame of $\OO_S(Y)$ defined by $s_Y=w_id_i$. Then
\[
 \frac{d_j}{d_i}=\frac{w_i}{w_j}
 =\frac{1-c_{ij}w_i}{g_{ij}}.
\]
Thus the frames $\ell_i=d_i\otimes\pi^*k_i$ of $L$ satisfy
\begin{equation}\label{eq:L-transition}
 \ell_j=(1-c_{ij}w_i)\ell_i.
\end{equation}
Their restrictions to $Y$ agree, in accordance with the trivialization
in \eqref{eq:restriction}.

\begin{proposition}\label{prop:no-metric}
The line bundle $L$ admits no $C^2$ Hermitian metric with semipositive
curvature.
\end{proposition}

\begin{proof}
Suppose such a metric exists, with $C^2$ psh weights $\varphi_i$ in the
frames $\ell_i$. By \eqref{eq:L-transition}, the functions
$\varphi_i|_Y$ agree on overlaps. They define a subharmonic function on
the compact connected curve $Y$, hence a constant. Rescaling the metric,
we may assume
\begin{equation}\label{eq:flat-value}
 \varphi_i(z,0)=0.
\end{equation}
In particular $\varphi_{i,z\bar z}(z,0)=0$. The Levi matrix at $(z,0)$
is positive semidefinite. A positive semidefinite Hermitian matrix with
a zero diagonal entry has zero entries in the corresponding row and
column. Therefore
\[
 \varphi_{i,w\bar z}(z,0)=0.
\]
The $C^1$ function
\[
 a_i(z)=\varphi_{i,w}(z,0)
\]
satisfies the Cauchy--Riemann equation and is holomorphic.

Equations \eqref{eq:metric-transition} and \eqref{eq:L-transition} give
\[
 \varphi_j(z,w_j)
 =\varphi_i(z,w_i)-\log|1-c_{ij}(z)w_i|^2,
\]
where the base coordinate change is understood. Differentiating with
respect to $w_i$ at zero and using \eqref{eq:normal-transition}, we obtain
\begin{equation}\label{eq:coboundary}
 g_{ij}a_j=a_i+c_{ij}.
\end{equation}
The holomorphic lifts $v_i'=v_i-a_ik_i$ consequently agree on overlaps:
\[
 v_j'=v_i+c_{ij}k_i-a_jg_{ij}k_i=v_i-a_ik_i=v_i'.
\]
They define a global holomorphic lift of $1$, splitting
\eqref{eq:extension-intro}. This is a contradiction.
\end{proof}

This calculation is the specialized first-order obstruction occurring
in \cite[Theorem~1.4 and Example~4.1]{Koike2021}. Related obstruction
classes, including their applicability to $C^2$ metrics, are discussed
in \cite{Wu2026}. Neither reference is needed for the preceding proof.

\subsection{A local continuous metric near a negative curve}\label{sec:cusp}

We isolate the analytic estimate used in the construction. The unitary
flatness assumption means that a smooth metric of zero curvature exists
on the restricted line bundle.

\begin{lemma}[Cusp correction]\label{lem:cusp}
Let $D$ be a compact connected smooth curve in a complex surface $X$. Suppose
$\deg N_{D/X}<0$, and let $A$ be a holomorphic line bundle such that
$A|_D$ is unitary flat. Choose a smooth metric $h_0$ on $A$ restricting
to a flat metric on $D$, and a smooth metric $h_D$ on $\OO_X(D)$ whose
curvature is strictly negative on $TD$. Put
\[
 q=\|s_D\|_{h_D}^2.
\]
For every $0<a\le\tfrac12$, there is a constant $\Lambda>0$ such that
the local weights of $h_0$ plus $\Lambda q^a$ are strictly psh off $D$
on a sufficiently small neighborhood of $D$ and extend as continuous
psh weights across $D$. They are locally of class $C^{0,2a}$; in
particular they are locally Lipschitz when $a=\tfrac12$.
For $0<a<\tfrac12$ one may take $\Lambda=1$ after shrinking the
neighborhood. At $a=\tfrac12$, a sufficiently large $\Lambda$ suffices.
\end{lemma}

\begin{proof}
The chosen metrics exist. Negative degree on the curve gives a negative
curvature metric on $N_{D/X}=\OO_X(D)|_D$. Any metric on a restricted
line bundle extends smoothly to a neighborhood: extend the logarithm
of its ratio to a fixed smooth metric. The same argument extends the
flat metric on $A|_D$.

Use finitely many relatively compact coordinates $(z,w)$ near $D$,
with $D=\{w=0\}$. Write
\[
 q=|w|^2e^{-b(z,w)},\qquad \beta=-b,
\]
where $b$ is the smooth weight of $h_D$. Negativity along $TD$ says
\begin{equation}\label{eq:beta-positive}
 \beta_{z\bar z}(z,0)>0.
\end{equation}
For $w\ne0$, set $\tau=\log q=\log|w|^2+\beta$. The chain rule gives
the exact identity
\begin{equation}\label{eq:cusp-exact}
 \Levi_{q^a}(v)
 =q^a\bigl(a\Levi_{\tau}(v)+a^2|\partial\tau(v)|^2\bigr).
\end{equation}
Evaluate this on
\[
 v=\xi\frac{\partial}{\partial z}
    +w\nu\frac{\partial}{\partial w}.
\]
Since $\log|w|^2$ is pluriharmonic off $D$,
\[
 \Levi_{\tau}(v)
 =\beta_{z\bar z}(z,0)|\xi|^2
   +O(|w|)(|\xi|^2+|\nu|^2),
\]
and
\[
 \partial\tau(v)
 =\nu+\beta_z(z,0)\xi+O(|w|)(|\xi|+|\nu|).
\]
Therefore
\begin{align}\label{eq:scaled-cusp}
 \frac{\Levi_{q^a}(\xi,w\nu)}{q^a}
 &=a\beta_{z\bar z}(z,0)|\xi|^2
   +a^2|\nu+\beta_z(z,0)\xi|^2\notag\\
 &\hspace{1em}+O(|w|)(|\xi|^2+|\nu|^2).
\end{align}
All remainder bounds are uniform after choosing the finite cover and
shrinking its coordinate neighborhoods.

The leading Hermitian form in \eqref{eq:scaled-cusp} is positive
definite: vanishing of its first term forces $\xi=0$ by
\eqref{eq:beta-positive}, and vanishing of the second then forces
$\nu=0$. Its smallest eigenvalue has a positive uniform lower bound on
the compact coordinate pieces covering $D$. Also $q^a$ is uniformly
comparable to $|w|^{2a}$. Thus, after shrinking the neighborhoods,
\begin{equation}\label{eq:cusp-lower}
 \Levi_{q^a}(\xi,w\nu)
 \ge c|w|^{2a}(|\xi|^2+|\nu|^2)
\end{equation}
for a uniform $c>0$.

Let $\varphi_0$ be a local weight of $h_0$. Its restriction to $D$ has
zero curvature, so $\varphi_{0,z\bar z}(z,0)=0$. Smoothness gives
\[
 \varphi_{0,z\bar z}=O(|w|),\qquad
 \varphi_{0,z\bar w}=O(1),\qquad
 \varphi_{0,w\bar w}=O(1).
\]
Consequently, in the same scaled variables,
\begin{equation}\label{eq:baseline-bound}
 |\Levi_{\varphi_0}(\xi,w\nu)|
 \le C|w|(|\xi|^2+|\nu|^2).
\end{equation}
Combining \eqref{eq:cusp-lower} and \eqref{eq:baseline-bound} yields
\begin{equation}\label{eq:cusp-coefficient}
 \Levi_{\varphi_0+\Lambda q^a}(\xi,w\nu)
 \ge\bigl(\Lambda c|w|^{2a}-C|w|\bigr)
       (|\xi|^2+|\nu|^2).
\end{equation}
If $0<a<\tfrac12$, take $\Lambda=1$ and then shrink until the
coefficient is positive. If $a=\tfrac12$, fix $\Lambda>C/c$; the
coefficient becomes $(\Lambda c-C)|w|>0$ throughout the punctured
neighborhood where the two uniform estimates hold. The constants
$c,C$ were obtained before choosing $\Lambda$, so this choice is not
circular. Every tangent vector at $w\ne0$ can be written in these
scaled variables. This proves strict plurisubharmonicity off $D$ in
both cases.

The function $q^a=|w|^{2a}e^{-ab}$ is continuous and locally
$C^{0,2a}$, including $2a=1$. The local weights
$\varphi_0+\Lambda q^a$ are bounded near $D$ and extend continuously
to it. The removable singularity theorem for psh functions across an
analytic hypersurface therefore proves that these extended weights
are psh. Adding the globally defined function $\Lambda q^a$ preserves
the metric transition law \eqref{eq:metric-transition}.
\end{proof}

\begin{remark}
The scale in \eqref{eq:scaled-cusp} is essential. The normal component
of the Levi form has order $|w|^{2a-2}$, while the tangential component
has order $|w|^{2a}$. These two scales allow the cusp to dominate the
bounded mixed curvature of a smooth initial metric. At $a=\tfrac12$
the scaled cusp and the initial error both have order $|w|$, which
explains the coefficient choice in \eqref{eq:cusp-coefficient}.
Ordinary smoothing of the cusp would not preserve this domination at
$D$. The argument proves no $C^1$ regularity.
\end{remark}

\subsection{An explicit global continuous metric}\label{sec:global}

Return to the ruled surface $S$ and the line bundle $L$ in
\eqref{eq:main-line}. Choose a positive smooth metric $h_K$ on $K$.
Using \eqref{eq:restriction}, extend the dual metric $h_K^{-1}$ on
$\OO_S(Y)|_Y$ to a smooth metric $h_Y$ on $\OO_S(Y)$. Set
\begin{equation}\label{eq:baseline-metric}
 h_0=h_Y\otimes\pi^*h_K,\qquad
 \theta_0=\Theta_{h_0}(L),\qquad
 q=\|s_Y\|_{h_Y}^2.
\end{equation}
The metric $h_0|_Y$ is the flat metric under the canonical
trivialization of $L|_Y$, and
\begin{equation}\label{eq:curvature-restrictions}
 \theta_0|_{TY}=0,\qquad
 \Theta_{h_Y}(\OO(Y))|_{TY}=-\Theta_{h_K}(K)<0.
\end{equation}

Fix $0<a\le\tfrac12$ and choose $\Lambda>0$ as in
Lemma~\ref{lem:cusp}. Since $q$ vanishes exactly on $Y$ and $S$ is
compact, there exists $\varepsilon>0$ such that the weights
$\varphi_0+\Lambda q^a$ are psh on $\{q<\varepsilon\}$.
Write $Q=\max_Sq$ and choose $0<\varepsilon<Q$. Choose a constant
\begin{equation}\label{eq:B-choice}
 B>\max_{\varepsilon\le t\le Q}\bigl(\Lambda t^a-\log t\bigr).
\end{equation}
Define, on $[0,Q]$,
\begin{equation}\label{eq:chi}
 \chi(0)=0,\qquad
 \chi(t)=\max\{\Lambda t^a,\log t+B\}\quad(t>0),
\end{equation}
and put
\begin{equation}\label{eq:global-metric}
 h=h_0\exp\bigl(-\chi(q)\bigr).
\end{equation}

\begin{proposition}\label{prop:continuous-metric}
The metric \eqref{eq:global-metric} is a positive continuous Hermitian
metric on $L$ with psh local weights. It is locally $C^{0,2a}$,
and is locally Lipschitz for $a=\tfrac12$.
\end{proposition}

\begin{proof}
The logarithmic branch in \eqref{eq:chi} tends to $-\infty$ as $t\to0$,
so $\chi(t)=\Lambda t^a$ for all sufficiently small $t$. Hence $\chi$ is
continuous at zero and bounded on $[0,Q]$.

On $S\setminus Y$, Poincar\'e--Lelong and
\eqref{eq:baseline-metric} give
\begin{equation}\label{eq:log-curvature}
 \theta_0+\ddc\log q=\pi^*\Theta_{h_K}(K)\ge0.
\end{equation}
As a current on $S$, the same expression equals
$\pi^*\Theta_{h_K}(K)+[Y]$. Thus $\varphi_0+\log q+B$ is a singular
psh weight of $L$.

On $\{0<q<\varepsilon\}$, the weights of $h$ are
\[
 \varphi_0+\chi(q)
 =\max\{\varphi_0+\Lambda q^a,\varphi_0+\log q+B\},
\]
a maximum of two psh functions. They are therefore psh. On
$\{q\ge\varepsilon\}$, \eqref{eq:B-choice} makes the logarithmic
branch strictly dominant. This domination also holds on a collar of
$\{q=\varepsilon\}$, so no boundary gluing issue arises. The weights
are psh there by \eqref{eq:log-curvature}. Near $Y$, the cusp branch is
dominant, and Lemma~\ref{lem:cusp} proves positivity across $Y$.

All the weights are finite and continuous. Near $Y$ their only
nonsmooth term is $\Lambda q^a=\Lambda|w|^{2a}e^{-ab}$,
which is $C^{0,2a}$ for $0<2a\le1$.
Away from $Y$, both branches are smooth, and their maximum is locally
Lipschitz, hence locally $C^{0,2a}$. This proves the asserted regularity.
\end{proof}

\begin{remark}[Bounded potentials and minimal singularities]\label{rem:minimal}
The current
\[
 T=\theta_0+\ddc\chi(q)\in c_1(L)
\]
has locally Lipschitz potentials when $a=\tfrac12$. It also has minimal
singularities: if $T'=\theta_0+\ddc v\ge0$ is another positive current
in $c_1(L)$, its quasi-psh potential $v$ is bounded above on compact
$S$, while $\chi(q)$ is bounded below. Hence
\[
 v\le\chi(q)+C_v
\]
for a constant $C_v$. This is precisely the comparison expressing that
$T$ is no more singular than $T'$. Here global potentials exist by the
$\partial\bar\partial$-lemma on the projective surface.

For this modified Grauert example, Proposition~\ref{prop:continuous-metric}
therefore gives an affirmative answer to the bounded-potential question
in \cite[Section~4]{Zhang2025} for the modified example. That question
also mentions Grauert's original example; the present argument does
not identify its line bundle with $L$ and does not settle that other
case. The metric is constructed here, not taken from that preprint.
This explicit answer for $L$ is not a statement about arbitrary nef
and big line bundles, nor a claim that every minimal-singularity
representative has continuous potentials.
\end{remark}

\section{Line-bundle exhaustions and the K\"ahler counterexample}\label{sec:radius}

\subsection{Dual norms and compact-group averaging}

\begin{lemma}\label{lem:norm}
Let $Z$ be a compact complex manifold and $A\to Z$ a holomorphic line
bundle with a positive continuous Hermitian metric $h$ having psh
local weights. Then the squared dual norm
\[
 \rho(v)=\|v\|_{h^{-1}}^2\qquad(v\in\Tot(A^{-1}))
\]
is a continuous psh exhaustion. If the weights are locally
$C^{0,\gamma}$ for $0<\gamma\le1$, so is $\rho$.
\end{lemma}

\begin{proof}
For a local frame $e_i$ of $A$, write a dual vector as
$v=\lambda_i e_i^*$, and write $h(e_i,e_i)=e^{-\varphi_i}$. Then
\begin{equation}\label{eq:local-norm}
 \rho=|\lambda_i|^2e^{\varphi_i}.
\end{equation}
Off the zero section, $\log\rho=\log|\lambda_i|^2+\varphi_i$ is psh.
Composition with the increasing convex function $\exp$ makes $\rho$
psh there. Formula \eqref{eq:local-norm} extends continuously by zero
to the zero section and is locally bounded, so removability gives psh
on the whole total space.

Fix a smooth Hermitian metric $h_{\mathrm{ref}}$ on $A^{-1}$. By
compactness of $Z$ and positivity of $h$, constants $c,C>0$ satisfy
\[
 c\|v\|_{h_{\mathrm{ref}}}^2
 \le\rho(v)\le C\|v\|_{h_{\mathrm{ref}}}^2.
\]
The closed disk bundle of any fixed radius over $Z$ is compact.
Every sublevel of $\rho$ is a closed subset of such a disk bundle and
is therefore compact. The H\"older assertion follows locally from
\eqref{eq:local-norm}.
\end{proof}

\begin{lemma}[Averaging preserves exhaustions]\label{lem:average}
Let a compact group $G$ act continuously on a locally compact Hausdorff
space $X$, and let $F$ be a continuous exhaustion. Its normalized Haar average
\[
 \overline F(x)=\int_G F(gx)\,dg
\]
is also an exhaustion. If $G$ is a compact Lie group, $X$ is a complex
manifold, the action is smooth and holomorphic, and $F$ is $C^2$ psh,
then $\overline F$ is $C^2$ psh. The same assertion holds in the
smooth category.
\end{lemma}

\begin{proof}
For $K_c=\{F\le c\}$, the orbit saturation $G\cdot K_c$ is compact.
If $x\notin G\cdot K_c$, then $F(gx)>c$ for every $g\in G$, and
continuity on the compact orbit implies $\overline F(x)>c$. Hence
\[
 \{\overline F\le c\}\subset G\cdot K_c.
\]
The sublevel on the left is closed and therefore compact. Differentiation
under the integral preserves $C^2$ regularity, and averaging the
nonnegative Levi forms of the holomorphic pullbacks preserves psh.
\end{proof}

\subsection{The radius criterion}

\begin{proposition}[The radius criterion]\label{prop:radius}
Let $Z$ be a compact complex manifold and $A\to Z$ a holomorphic line
bundle. If $\Tot(A^{-1})$ admits a $C^2$ psh exhaustion, then $A$ admits
a $C^2$ Hermitian metric with semipositive curvature. If the exhaustion
is smooth, the metric is smooth.
\end{proposition}

\begin{proof}
Average the exhaustion over the holomorphic circle action
$v\mapsto e^{i\alpha}v$ on the fibers, using
Lemma~\ref{lem:average}. Denote the resulting invariant exhaustion by
$F$. For a local frame $e_i$ of $A$ and $v=\lambda_i e_i^*$, circle
invariance gives
\begin{equation}\label{eq:radial-function}
 F(z,\lambda_i)=U_i(z,t),\qquad t=\log|\lambda_i|^2,
 \quad\lambda_i\ne0.
\end{equation}
The function
\begin{equation}\label{eq:holomorphic-pullback}
 V_i(z,\zeta)=U_i(z,\operatorname{Re}\zeta)
             =F(z,e^{\zeta/2})
\end{equation}
is a holomorphic pullback of $F$ and is psh. In particular
$U_{i,tt}\ge0$.

For fixed $z$, continuity at the zero section gives the finite limit
\[
 \lim_{t\to-\infty}U_i(z,t)=F(z,0).
\]
Exhaustiveness gives $U_i(z,t)\to+\infty$ as $t\to+\infty$.
A differentiable convex function with finite left limit has
nonnegative derivative. Moreover, if its derivative vanishes at
$t_0$, convexity forces it to be constant on $(-\infty,t_0]$.
Consequently
\begin{equation}\label{eq:radial-positive}
 U_{i,t}(z,t)>0\quad\text{whenever }U_i(z,t)>F(z,0).
\end{equation}

Fix $c>\max_ZF|_Z$. By convexity, the two limits, and
\eqref{eq:radial-positive}, there is a unique solution
\begin{equation}\label{eq:implicit-radius}
 U_i(z,\sigma_i(z))=c,
 \qquad U_{i,t}(z,\sigma_i(z))>0.
\end{equation}
The implicit function theorem makes $\sigma_i$ of class $C^2$, or
smooth when $F$ is smooth.

We verify the sign of its Levi form explicitly. Differentiating
\eqref{eq:implicit-radius} twice in base coordinates gives
\begin{align}\label{eq:implicit-hessian}
 0={}&U_{i,z_j\bar z_k}
 +U_{i,z_jt}\sigma_{i,\bar z_k}
 +U_{i,t\bar z_k}\sigma_{i,z_j}\notag\\
 &+U_{i,tt}\sigma_{i,z_j}\sigma_{i,\bar z_k}
 +U_{i,t}\sigma_{i,z_j\bar z_k},
\end{align}
with all derivatives of $U_i$ evaluated at $(z,\sigma_i(z))$.
For $V_i$ in \eqref{eq:holomorphic-pullback}, its $\zeta\bar\zeta$
entry is $U_{i,tt}/4$, and its mixed entries have a factor $1/2$.
Thus \eqref{eq:implicit-hessian} is precisely the identity
\begin{equation}\label{eq:radius-Levi}
 \Levi_{V_i}\bigl(\xi,2\partial\sigma_i(\xi)\bigr)
 =-U_{i,t}\Levi_{\sigma_i}(\xi).
\end{equation}
The left side is nonnegative and $U_{i,t}>0$, so
$\psi_i=-\sigma_i$ is psh.

If $e_j=g_{ij}e_i$, the dual coordinates satisfy
$\lambda_j=g_{ij}\lambda_i$. The level in
\eqref{eq:implicit-radius} consequently transforms by
\[
 \sigma_j=\sigma_i+\log|g_{ij}|^2,
\]
so
\[
 \psi_j=\psi_i-\log|g_{ij}|^2.
\]
These are exactly the metric transition laws
\eqref{eq:metric-transition}. Hence $\|e_i\|^2=e^{-\psi_i}$ defines
the desired semipositive Hermitian metric on $A$.
\end{proof}

\begin{corollary}\label{cor:criterion}
For a holomorphic line bundle $A$ over a compact complex manifold,
$\Tot(A^{-1})$ admits a $C^2$ psh exhaustion if and only if $A$ admits
a $C^2$ semipositive Hermitian metric. The same equivalence holds with
$C^2$ replaced by $C^\infty$.
\end{corollary}

\begin{proof}
The forward implication is Proposition~\ref{prop:radius}. For the
reverse implication, the squared dual norm in Lemma~\ref{lem:norm}
has the regularity of the given metric, by its local expression
\eqref{eq:local-norm}.
\end{proof}

\subsection{Proof of the K\"ahler counterexample}\label{sec:conclusion}

\begin{proof}[Proof of Theorem~\ref{thm:main}]
The construction in Subsection~\ref{sec:ruled} gives a smooth projective
surface $S$ and an algebraic line bundle $L$. The total space $M$ is
the open subset
\begin{equation}\label{eq:compactification}
 M=\PP_{\mathrm{lines}}(\OO_S\oplus L^{-1})
       \setminus\PP_{\mathrm{lines}}(L^{-1}).
\end{equation}
The identification sends $v\in L^{-1}_z$ to the line $[(1,v)]$.
The projective bundle in \eqref{eq:compactification} is a smooth
projective threefold. Restricting any of its K\"ahler forms proves
that $M$ is K\"ahler.

Let $h$ be the continuous semipositive metric from
Proposition~\ref{prop:continuous-metric}. If $p:M\to S$ is the bundle
projection, the exhaustion is explicitly
\begin{equation}\label{eq:explicit-exhaustion}
 \rho(v)=\|v\|_{h^{-1}}^2
       =\|v\|_{h_0^{-1}}^2
          \exp\!\bigl(\chi(q(p(v)))\bigr).
\end{equation}
Lemma~\ref{lem:norm} shows that $\rho$ is continuous, psh, and
exhaustive. Taking $a=\tfrac12$ and the coefficient $\Lambda$ supplied
by Lemma~\ref{lem:cusp} makes $h$, its weights, and $\rho$ locally
Lipschitz. In particular the same example has local $C^{0,\gamma}$
regularity for all $0<\gamma\le1$.

If a $C^2$ psh exhaustion on $M$ existed,
Proposition~\ref{prop:radius} would produce a $C^2$ semipositive
Hermitian metric on $L$. Proposition~\ref{prop:no-metric} excludes
such a metric. Therefore no $C^2$ psh exhaustion exists, completing
the proof.
\end{proof}

\begin{remark}[Interpretation]
The example separates existence of a continuous exhaustion from
existence of any smooth exhaustion. The conclusion does not depend on
whether the particular function \eqref{eq:explicit-exhaustion} can be
smoothed: every hypothetical $C^2$ exhaustion leads to the forbidden
metric. The argument does not determine whether a $C^1$ psh
exhaustion exists. Corollary~\ref{cor:quasiprojective-regularization}
shows that dimension three is minimal within smooth quasi-projective
manifolds; minimality among arbitrary noncompact K\"ahler manifolds
is not established. Thus a definition of weak pseudoconvexity
requiring a smooth
psh exhaustion is strictly stronger than the corresponding definition
requiring only a continuous psh exhaustion, even for smooth
quasi-projective K\"ahler threefolds.
\end{remark}

\section{Compact modifications and non-K\"ahler counterexamples}
\label{sec:nonkahler}

The K\"ahler counterexample can be changed over a single point to obtain
a non-K\"ahler counterexample of the same dimension.  We first isolate
the exhaustion argument, then insert Hironaka's opposite-order blowup
configuration in an exceptional plane.  Compact products give a simpler
alternative when an increase in dimension is allowed.

\subsection{Exhaustions under compactly supported modifications}

\begin{proposition}\label{prop:compact-modification}
Let $f:X'\to X$ be a proper surjective holomorphic map between nonempty
complex manifolds.  Suppose that, for a compact set $K\subset X$, it
restricts to a biholomorphism
\[
 X'\setminus f^{-1}(K)\longrightarrow X\setminus K.
\]
For every $r\in\{0,2,\infty\}$, $X$ admits a $C^r$ psh exhaustion
if and only if $X'$ does.
\end{proposition}

\begin{proof}
The case $K=\varnothing$ is immediate.  Holomorphic pullback preserves
plurisubharmonicity and the stated
regularity, and properness makes the sublevels of a pulled-back
exhaustion compact.  Conversely, suppose $F$ is a $C^r$ psh exhaustion
on $X'$.  Choose a relatively compact open neighborhood $U$ of $K$ and
\[
 c>\max_{f^{-1}(\overline U)}F.
\]
Choose a smooth convex nondecreasing function
$\chi:\R\to[0,\infty)$ that vanishes on $(-\infty,c]$ and tends to
infinity.  For example, integrate a smooth nondecreasing function
which is zero up to $c$ and one from $c+1$ onward.  On $X\setminus K$
define
\[
 G=\chi\circ F\circ(f|_{X'\setminus f^{-1}(K)})^{-1}.
\]
This function is $C^r$ and psh.  It vanishes on $U\setminus K$, so
extension by zero on $K$ gives a $C^r$ psh function on all of $X$.
For each nonempty sublevel $\{G\le a\}$ choose $b$ such that
$\chi(t)\le a$ implies $t\le b$.  Then
\[
 \{G\le a\}\subset K\cup f(\{F\le b\}).
\]
The left side is closed and the right side is compact.  Thus $G$ is an
exhaustion; no smooth descent of the original function $F$ is needed.
\end{proof}

\subsection{A non-K\"ahler modification over one point}

\begin{proposition}\label{prop:point-modification}
Let $X$ be a connected smooth complex threefold and $x\in X$.
There is a connected smooth non-K\"ahler threefold $N$ and a proper
holomorphic modification $\mu:N\to X$ inducing a biholomorphism
\[
 N\setminus\mu^{-1}(x)\simeq X\setminus\{x\}.
\]
\end{proposition}

\begin{proof}
Blow up $x$, obtaining
\[
 b:W=\Bl_xX\longrightarrow X,\qquad E=b^{-1}(x)\simeq\PP^2.
\]
In homogeneous coordinates $[u:v:w]$ on $E$, take
\[
 A=\{w=0\},\qquad B=\{uv=w^2\}.
\]
The line $A$ and smooth conic $B$ meet transversely at exactly
$p=[1:0:0]$ and $q=[0:1:0]$.  Near either point they are the two
coordinate axes in a smooth surface $E$ inside $W$.

Over $W\setminus\{q\}$, first blow up $A\setminus\{q\}$ and then
the strict transform of $B\setminus\{q\}$.  Over
$W\setminus\{p\}$, first blow up $B\setminus\{p\}$ and then the
strict transform of $A\setminus\{p\}$.  All these centers are closed
smooth curves in the respective open manifolds.  On
$W\setminus\{p,q\}$ the centers are disjoint, so the blowups commute
canonically over the base.  Glue the two resulting manifolds through
this isomorphism and write
\[
 \nu:N\longrightarrow W,\qquad \mu=b\circ\nu.
\]
The glued space is Hausdorff: points with distinct images can be
separated in $W$, and points with the same image belong to one of the
Hausdorff open manifolds lying over the two base opens.  It is smooth
of dimension three because every blowup center is smooth.  Properness
is local on the target, so $\nu$ and $\mu$ are proper.  Each blowup
chart is connected and the charts overlap, hence $N$ is connected.
All centers lie in $E$, so $\mu$ is an isomorphism away from $x$.
This is Hironaka's opposite-order construction
\cite[Section~6]{Barlet2015}; see also \cite[Tag~08J1]{StacksHironaka}.

For completeness, the obstruction to a K\"ahler form is entirely
contained in compact exceptional surfaces.  Let $D_A,D_B\subset N$
be the exceptional divisors dominating $A,B$, and let
$\alpha,\beta\in H_2(N,\R)$ be the classes of their general
$\PP^1$ fibers.  The surfaces $D_A,D_B$ are smooth and compact:
locally each is a ruled surface or its blowup at a point, and they
map properly to the compact curves $A,B$.  The fiber divisors of
each map are linearly equivalent, since $A\simeq B\simeq\PP^1$,
and hence are homologous on the corresponding compact surface.

Over $p$, the strict transform of $B$ meets the first exceptional
divisor transversely at one point.  The second blowup therefore
restricts on that divisor to a surface blowup at that point.  Its
fiber over $p$ is the reduced union of the strict transform
$R_p\simeq\PP^1$ of the first exceptional fiber and the new
exceptional fiber $T_p\simeq\PP^1$. Both occur with multiplicity one.
Indeed, take local coordinates $(x,y,z)$ with
$A=\{y=z=0\}$ and $B=\{x=z=0\}$. In the first blowup chart
$z=yt$, the first exceptional divisor is $y=0$ and the second center
is $x=t=0$. Its restriction to $y=0$ is the blowup of the smooth
surface with coordinates $(x,t)$ at the origin. The pullback of the
fiber equation $x=0$ is its strict transform plus the exceptional
curve, each once. The curve $T_p$ is also the fiber of $D_B\to B$
over $p$. Thus
\[
 \alpha=[R_p]+[T_p],\qquad \beta=[T_p].
\]
Over $q$ the order is reversed.  With analogous notation,
\[
 \beta=[R_q]+[T_q],\qquad \alpha=[T_q].
\]
These relations hold in ordinary homology: they arise from fiber
divisors on compact $D_A$ and $D_B$, regardless of compactness of $N$.
Adding them gives
\begin{equation}\label{eq:hironaka-zero-cycle}
 [R_p]+[R_q]=0\quad\text{in }H_2(N,\R).
\end{equation}
If $\omega$ were a K\"ahler form on $N$, closedness and
\eqref{eq:hironaka-zero-cycle} would imply
\[
 0=\int_{R_p}\omega+\int_{R_q}\omega>0,
\]
a contradiction.
\end{proof}

\begin{proof}[Proof of Theorem~\ref{thm:nonkahler-threefold}]
Apply Proposition~\ref{prop:point-modification} to the threefold $M$
of Theorem~\ref{thm:main} at the chosen point $x$.  If $\rho$ is its
continuous psh exhaustion, then
\[
 \rho_N=\rho\circ\mu,\qquad
 \{\rho_N\le a\}=\mu^{-1}(\{\rho\le a\})
\]
is a continuous psh exhaustion of $N$.  A holomorphic map is locally
Lipschitz in smooth coordinate charts, so local Lipschitz regularity
of the exhaustion is preserved.  A $C^2$ psh
exhaustion on $N$ would give one on $M$ by
Proposition~\ref{prop:compact-modification} with $K=\{x\}$, contrary
to Theorem~\ref{thm:main}.
\end{proof}

\begin{remark}[The sense in which the change is local]
The point $x$ can be chosen outside the zero section of $M\to S$.
The construction then preserves a neighborhood of that section and,
by properness, a neighborhood of infinity.  More precisely, for every
neighborhood $U$ of $x$, $\mu$ is an isomorphism over $M\setminus U$.
This is a modification supported over one point; it makes no claim
about a small deformation of the complex structure on a fixed smooth
manifold.
\end{remark}

\subsection{Compact products and an explicit fivefold}\label{sec:products}

\begin{proposition}[Compact product principle]\label{prop:compact-product}
Let $X$ be a nonempty complex manifold and $T$ a nonempty compact
connected complex manifold.  For each $r\in\{0,2,\infty\}$, $X$
admits a $C^r$ psh exhaustion if and only if $X\times T$ does.  If
$T$ is non-K\"ahler, then $X\times T$ is non-K\"ahler.
\end{proposition}

\begin{proof}
Pullback by the projection preserves regularity and
plurisubharmonicity, with compact sublevels
\begin{equation}\label{eq:product-sublevels}
 \{f\circ\operatorname{pr}_X\le a\}=\{f\le a\}\times T.
\end{equation}
Conversely, restriction to the closed holomorphic slice
$X\times\{t_0\}$ preserves regularity, plurisubharmonicity, and
compact sublevels.  A K\"ahler form on the product would restrict to
one on $\{x_0\}\times T$, proving the final assertion.  In fact, the
maximum principle makes every continuous psh function on the product
constant on each compact connected fiber $\{x\}\times T$.
\end{proof}

\begin{corollary}[A non-K\"ahler fivefold]\label{cor:nonkahler}
Let $M$ be the threefold of Theorem~\ref{thm:main} and let
\begin{equation}\label{eq:hopf-surface}
 H=(\C^2\setminus\{0\})/\langle z\mapsto2z\rangle.
\end{equation}
Then $M\times H$ is a smooth non-K\"ahler fivefold admitting a
continuous psh exhaustion but no $C^2$ psh exhaustion.  For every
$0<\gamma<1$ its exhaustion can be chosen locally $C^{0,\gamma}$.
\end{corollary}

\begin{proof}
The dilation action is free and properly discontinuous: if a compact
set lies in $r\le\|z\|\le R$, only finitely many integers $n$ satisfy
$r/R\le2^n\le R/r$.  Polar coordinates identify its quotient smoothly
with
\begin{equation}\label{eq:hopf-topology}
 H\simeq S^3\times\bigl(\R/(\log2)\mathbb Z\bigr)
 \simeq S^3\times S^1.
\end{equation}
Thus $H$ is a compact connected complex surface with
$H^2_{\mathrm{dR}}(H)=0$.  Any K\"ahler form $\omega$ would be exact,
and Stokes' theorem would give $\int_H\omega^2=0$, contradicting
positivity.  Apply Proposition~\ref{prop:compact-product} and
Theorem~\ref{thm:main}.  Explicitly,
\begin{equation}\label{eq:nonkahler-exhaustion}
 \rho_{M\times H}(v,[z])=\rho(v)
\end{equation}
is the required exhaustion, with the same local H\"older regularity.
\end{proof}

\section{Surface regularization and boundary existence criteria}\label{sec:surfaces}
We prove that a continuous psh exhaustion on a divisor complement in
a compact K\"ahler surface implies a smooth one, and hence obtain
regularization for every smooth complex quasi-projective surface.
The averaging and current argument below treats every reducible or
disconnected boundary and requires no minimal-model reduction.
We distinguish this conditional regularization theorem from the
separate numerical and Ueda-theoretic criteria that construct an
exhaustion without assuming a continuous one.

Unless otherwise specified, a divisor compactification means
$X=S\setminus\supp D$, where $S$ is a smooth compact connected
complex surface and $D\ne0$ is effective. The Enriques--Kodaira
classification concerns $S$, not $X$. Table~\ref{tab:surface-status}
summarizes the consequences for such compactifiers. No divisorial
compactification is assumed to exist for an arbitrary noncompact
surface.

\subsection{Two basic affirmative cases}\label{sec:surface-curves}

\begin{proposition}[A known holomorphic-convexity consequence]
\label{prop:surface-nonconstant-function}
Let $X$ be a connected noncompact complex manifold of dimension two
with a continuous psh exhaustion. If $\OO(X)\ne\C$, then $X$ has a smooth psh exhaustion.
\end{proposition}
\begin{proof}
The holomorphic-convexity theorem of V\^{\i}j\^{\i}itu
\cite[Theorem~1.1]{Vajaitu2022} makes $X$ holomorphically convex.
Its Remmert reduction $r:X\to Y$ is proper with Stein target. A smooth
strongly psh exhaustion of the complex space $Y$ pulls back to a smooth
psh exhaustion of $X$. Here smoothness on $Y$ means local restriction
of smooth ambient functions, so the pullback is smooth even over
singular points of $Y$.
\end{proof}

This is a consequence of the cited theorem, not a new
holomorphic-convexity result. This proposition alone leaves the case $\OO(X)=\C$
untreated; Theorem~\ref{thm:divisor-regularization} below also covers
that case under its compactification hypothesis.
Likewise, the regularization theorem of Mongodi--Slodkowski for domains
in a surface already carrying a real-analytic psh exhaustion
\cite[Corollaries~4.4--4.5]{MongodiSlodkowski2020} requires that ambient
hypothesis; one cannot assume it for $X$ itself.

\begin{proposition}[Line-bundle total spaces over a curve]
\label{cor:curve-total-space}
Let $A$ be a holomorphic line bundle on a compact connected Riemann
surface $B$. The following are equivalent:
\[
 \begin{gathered}
 \Tot(A^{-1})\text{ has a continuous psh exhaustion}\\
 \Longleftrightarrow\quad \deg A\ge0\quad\Longleftrightarrow\\
 \Tot(A^{-1})\text{ has a smooth psh exhaustion}.
 \end{gathered}
\]
\end{proposition}
\begin{proof}
Choose a smooth reference metric with curvature $\theta_0$ and an
area form $\omega$ with $\int_B\omega=1$. For $d=\deg A\ge0$, solve
$\ddc f=d\omega-\theta_0$. The resulting metric has curvature
$d\omega$, and its squared dual norm is a smooth psh exhaustion by
Lemma~\ref{lem:norm}. Conversely, the compactification
\[
 \Tot(A^{-1})=\PP_{\mathrm{lines}}(\OO_B\oplus A^{-1})
                 \setminus\PP_{\mathrm{lines}}(A^{-1})
\]
has infinity section with normal bundle $A$. If $\deg A<0$, that
section has negative self-intersection, contradicting the boundary
obstruction in Lemma~\ref{lem:surface-boundary-tools}\textup{(i)} below.
\end{proof}

Thus the line-bundle construction of Theorem~\ref{thm:main} cannot be
lowered to dimension two by replacing its compact base by a curve.
The full quasi-projective surface case is covered by
Corollary~\ref{cor:quasiprojective-regularization}.

\subsection{Boundary tools and permitted modifications}
\label{sec:surface-pair-reduction}

For a divisor complement, a continuous exhaustion tends uniformly to
$+\infty$ at the boundary: each sublevel is a compact subset of $X$,
hence its complement in $S$ is a neighborhood of $\supp D$. We call a
psh function on a punctured boundary neighborhood which diverges in
this way a \emph{divergent collar}.

Blowing up a boundary point preserves $X$ only when the new boundary
is the \emph{full reduced inverse image}, including the exceptional
curve. This permits resolution of the boundary without changing $X$.
Blowups at interior points give proper modifications with compact
exceptional locus, and Proposition~\ref{prop:compact-modification}
preserves both exhaustion properties. By contrast, contracting a curve
which meets but is not contained in the boundary need not induce a
proper map of the complements. We make no unrestricted reduction of
projective pairs to minimal compactifiers.

\begin{lemma}[Boundary obstruction and collar construction]
\label{lem:surface-boundary-tools}
Let $S$ be smooth compact and let $D\ne0$ be effective.
\begin{enumerate}[label=\textup{(\roman*)},leftmargin=2.5em]
\item\label{item:negative-boundary}
If a connected component of $\supp D$ has negative definite
intersection matrix, then $S\setminus\supp D$ has no continuous psh
exhaustion.
\item\label{item:collar-gluing}
If every connected boundary component admits a smooth divergent psh
collar, then $S\setminus\supp D$ has a smooth psh exhaustion.
\item\label{item:matrix-alternatives}
Let $E=\bigcup E_i$ be a connected reduced simple normal crossing
\textup{(SNC)} divisor and let $\lambda$ be the largest eigenvalue of
$I=(E_i\cdot E_j)$. If $\lambda>0$, $E$ admits a smooth strictly psh
divergent collar. If $\lambda=0$, the matrix is negative semidefinite
with one-dimensional kernel, and there are positive integers $a_i$
with
\begin{equation}\label{eq:surface-null-boundary}
 A=\sum_i a_iE_i,\qquad A\cdot E_i=0,\qquad A^2=0.
\end{equation}
When $S$ is projective, this divisor $A$ is nef.
\end{enumerate}
\end{lemma}
\begin{proof}
For \textup{(i)}, Grauert's contraction theorem \cite[\S4, subsection~8(e), pp.~366--367]{Grauert1962}
contracts the indicated component to a point $p$ of a normal surface,
without changing a punctured neighborhood. A continuous psh function
on such a punctured neighborhood is bounded above near $p$. Indeed,
a local finite projection to a bidisc, with $p$ the only point over
$0$, sends it to the fiberwise maximum $v_*$, a continuous psh function
on the punctured bidisc. This maximum is psh off the branch locus;
continuity and local boundedness extend this across the branch locus
away from $0$. On
\[
 \{|z_1|\le r,\ |z_2|=r\}
\]
it has a finite maximum. The maximum principle on each vertical disc
with $z_1\ne0$, followed by continuity at $z_1=0$, gives the same bound
on the punctured smaller bidisc. This contradicts divergence of an
exhaustion at the contracted boundary.

For \textup{(ii)}, choose disjoint smaller boundary neighborhoods $W_E$
compactly contained in the collar neighborhoods and choose $T_E$ larger
than the collar function on $\partial W_E$. A smooth convex
nondecreasing function $\chi_E$, zero on $(-\infty,T_E]$ and tending to
infinity, makes $\chi_E(\Phi_E)$ vanish near $\partial W_E$. Extend by
zero outside $W_E$ and sum over the finitely many boundary components.
The result is smooth and psh, diverges at every boundary point, and is
proper because $S$ is compact.

For \textup{(iii)}, the off-diagonal entries of $I$ are nonnegative and
its graph is connected. Perron--Frobenius applied after a scalar shift
makes $\lambda$ simple with a strictly positive eigenvector. If
$\lambda>0$, rational approximation and clearing denominators give
$A=\sum a_iE_i$ with $a_i>0$ and $A\cdot E_i>0$ for every $i$.
We spell out the metric construction needed here. The restriction
$\OO_S(A)|_E$ is ample. Embed $E$ by a power of this line bundle and
use its Fubini--Study metric. Lift the embedding sections locally to
$S$ and take the appropriate metric root. Their ambient extensions
have positive tangential curvature along the smooth part of $E$;
at a node they have positive curvature in all directions, since the
Zariski tangent space of $E$ there is the ambient tangent space and
the embedding is injective on tangent spaces. Patch these extensions
using a partition of unity relative to a reference metric, keeping
the chosen metrics near the nodes. The restriction to $E$, and hence
the tangential curvature on its smooth part, is unchanged.

Let $q=\|s_E\|^2$ for any smooth metric on $\OO_S(E)$. On the smooth
part of $E$, $\ddc q$ is positive in the normal direction and zero
in the tangent direction; at a node it is zero. Adding $Cq$ to the
patched weights, for sufficiently large $C$, makes the curvature
strictly positive along all of $E$, by the Schur complement and
compactness, and hence in a neighborhood. For this metric $h$,
$-\log\|s_A\|_h^2$ is the required collar.

If $\lambda=0$, simplicity gives a one-dimensional positive kernel,
which has a rational generator because $I$ is integral. Clearing
denominators proves \eqref{eq:surface-null-boundary}. Intersections
with curves outside $E$ are nonnegative by effectiveness; in the
projective case this proves nefness.
\end{proof}

In particular, under the continuous-exhaustion hypothesis a connected
SNC boundary is either already handled by a positive collar, or
supports the weighted square-zero divisor
\eqref{eq:surface-null-boundary}. A negative definite boundary is an
\emph{exclusion}, not a continuous-versus-smooth counterexample.

\subsection{Regularization for divisor complements in compact K\"ahler surfaces}
\label{subsec:divisor-regularization}
The continuous-to-smooth implication for divisor complements follows
from positive $dd^c$-closed currents on compact K\"ahler surfaces.
The argument applies to reducible, disconnected, and singular divisors
and does not require a minimal compactifier. We use the normalization
of Section~\ref{sec:conventions}. A modification of a Stein space
means a proper holomorphic modification onto a Stein space.

\begin{theorem}[Regularization for compact K\"ahler divisor complements]
\label{thm:divisor-regularization}
Let $S$ be a smooth compact connected K\"ahler surface and let
$D\ne0$ be a reduced effective divisor.  If $X=S\setminus D$ admits
a continuous plurisubharmonic exhaustion, then it admits a smooth
plurisubharmonic exhaustion.

More precisely, if $X$ is not a modification of a Stein space, there
are positive integers $a_i$, one for every irreducible component $D_i$
of $D$, such that
\[
 A=\sum_i a_iD_i,\qquad A\cdot D_i=0\ \text{for every }i,
 \qquad A^2=0,
\]
and $\mathcal O_S(A)$ admits a smooth Hermitian metric with zero
curvature on a neighborhood of $D$.  In particular, a punctured
neighborhood of the entire boundary admits a smooth pluriharmonic
function tending to $+\infty$ along $D$.
\end{theorem}

We first prove the averaging statement that supplies the current used
in the proof.  Its conclusion concerns Lelong numbers at every point,
including points at which a weak limit could acquire new concentration.

\begin{lemma}[Currents obtained by averaging exhaustion levels]
\label{lem:level-averaging}
Let $(S,\omega)$ be a compact K\"ahler surface, and let $X\subset S$
be a connected noncompact open subset with a continuous plurisubharmonic
exhaustion $\rho$.  If $X$ is not a modification of a Stein space,
there is a positive $(1,1)$-current $T$ on $S$ such that
\begin{equation}
\label{eq:averaged-current-properties}
 \ddc T=0,\qquad
 \int_S T\wedge\omega=1,\qquad
 \operatorname{supp}T\Subset X,\qquad
 \nu(T,x)=0\quad(x\in S).
\end{equation}
\end{lemma}

\begin{proof}
Let $\Sigma_X^0$ be the continuous minimal kernel: the intersection,
over all continuous plurisubharmonic exhaustions $\psi$ of $X$, of
the sets where $\psi$ is not strongly plurisubharmonic.  Here strong
plurisubharmonicity means local stability under sufficiently small
$C^2$ perturbations.  By \cite[Lemma~2.1]{MongodiSlodkowski2020}, if
$\rho(\Sigma_X^0)$ contains no upper half-line, then $X$ is a
modification of a Stein space.  Thus, under the present hypothesis,
there is $t_0\in\mathbb R$ such that
\begin{equation}
\label{eq:kernel-halfline}
 (t_0,+\infty)\subset\rho(\Sigma_X^0).
\end{equation}
This invocation uses the continuous version of that lemma, which is
stated for arbitrary complex manifolds.

Fix $t>t_0$ and put $K_t=\{\rho=t\}$.  We claim that no smooth
strictly plurisubharmonic function exists on a neighborhood of the
whole compact set $K_t$.  Suppose, to the contrary, that $v$ is such
a function on a neighborhood $W$.  Properness of $\rho$ allows us
to choose $\delta>0$ with
\[
 B=\{t-\delta\leq\rho\leq t+\delta\}\Subset W.
\]
After a positive affine change of $v$, assume $0<v<1$ on $B$.
Choose $0<\eta<1$, $0<\varepsilon<\eta\delta/2$, and set
$g(s)=s+\eta|s-t|$.  The function $g$ is convex, increasing, and
tends to $+\infty$ as $s\to+\infty$.  Define on $X$
\[
 \psi(x)=
 \begin{cases}
  \max\{g(\rho(x)),\,\rho(x)+\varepsilon v(x)\},
       & |\rho(x)-t|<\delta,\\
  g(\rho(x)),& |\rho(x)-t|\geq\delta.
 \end{cases}
\]
On the part of the band where $|\rho-t|>\delta/2$, the first branch
strictly dominates, so the definitions glue.  On $K_t$ the second
branch strictly dominates, and hence it also dominates on a
neighborhood of $K_t$.  Consequently $\psi$ is a continuous
plurisubharmonic exhaustion strongly plurisubharmonic near $K_t$.
This contradicts $K_t\cap\Sigma_X^0\ne\varnothing$, which follows
from \eqref{eq:kernel-halfline}.

We now apply Hahn--Banach separation on $K_t$.  Let $\mathcal E_t$
be the real Banach space of continuous sections over $K_t$ of the
bundle of real $(1,1)$-forms on $S$.  Its open convex cone
$\mathcal P_t$ consists of forms that are positive definite at every
point of $K_t$.  The linear subspace
\[
 \mathcal L_t
   =\{\ddc f|_{K_t}:f\in C^\infty(S,\mathbb R)\}
\]
does not meet $\mathcal P_t$, by the preceding claim.  Conversely,
a smooth function defined near $K_t$ can be extended to a smooth
function on $S$ agreeing with it near $K_t$, by a cutoff; thus there
is no distinction between the local and global formulations here.
Separating the linear subspace from the open cone gives a nonzero
continuous real linear functional $\Lambda_t$ vanishing on
$\mathcal L_t$ and nonnegative on $\mathcal P_t$, and hence on its
closure.  The subspace need not be closed for this separation.
Restriction of test forms to $K_t$ turns $\Lambda_t$ into an
order-zero positive $(1,1)$-current supported on $K_t$.  Its
annihilation of $\ddc f$ for all smooth $f$ says that it is
$dd^c$-closed.  Its pairing with $\omega$ is strictly positive:
otherwise positivity and the inequalities $-c\omega\leq\gamma
\leq c\omega$ on $K_t$ would make it vanish on every
$\gamma\in\mathcal E_t$.  We may therefore normalize it to obtain
\begin{equation}
\label{eq:level-current}
 T_t\geq0,\qquad \ddc T_t=0,\qquad
 \operatorname{supp}T_t\subset K_t,\qquad
 \int_S T_t\wedge\omega=1.
\end{equation}

Fix $t_0<a<b$.  For $n\geq1$ choose the equally spaced levels
\[
 t_{n,j}=a+\frac{j-\tfrac12}{n}(b-a),\qquad 1\leq j\leq n,
 \qquad T_n=\frac1n\sum_{j=1}^n T_{t_{n,j}}.
\]
Only finitely many arbitrary choices are needed for each $n$;
no measurable selection is involved.  Every $T_n$ has mass one and
is supported in the fixed compact shell
$K=\rho^{-1}([a,b])\Subset X$.  Weak compactness of positive
currents of bounded mass on $S$ yields a subsequence, still denoted
by $T_n$, converging to a positive current $T$.  Testing against
$\omega$ and against $\ddc f$ shows that $T$ has mass one and is
$dd^c$-closed.  Testing on $S\setminus K$ shows
$\operatorname{supp}T\subset K$.

It remains to rule out Lelong concentration.  Fix $x\in X$ and
choose holomorphic coordinates on a neighborhood $U\Subset X$
containing $\overline{B(x,2R)}$.  Write
$\beta=\ddc|z|^2$.  The Lelong--Jensen monotonicity for positive
$dd^c$-closed currents, recalled in
\cite[Appendix~B, equations~(B.1)--(B.2)]{DNS22}, implies that for
every current in \eqref{eq:level-current} and $0<r<R$,
\begin{equation}
\label{eq:uniform-lelong-bound}
 \frac{1}{r^2}\int_{B(x,r)}T_t\wedge\beta
 \leq\frac{1}{R^2}\int_{B(x,R)}T_t\wedge\beta
 \leq C_R.
\end{equation}
Indeed, $\beta\leq c\omega$ on $\overline{B(x,R)}$, so one may
take $C_R=c/R^2$ using the mass normalization.  Thus $C_R$ is
independent of $t$, $n$, and $j$.  Monotonicity applies to open
balls as well; alternatively it follows for them by increasing
limits from smaller radii.  The positive factor between the
$dd^c$ convention in \cite{DNS22} and ours changes neither this
monotonicity nor vanishing of Lelong numbers.

If $T_{t_{n,j}}$ contributes mass to $B(x,r)$, then its supporting
level belongs to
\[
 I_r=\left[\inf_{B(x,r)}\rho,\ \sup_{B(x,r)}\rho\right].
\]
The number of grid points in $I_r$ is at most
$n|I_r|/(b-a)+2$.  Hence \eqref{eq:uniform-lelong-bound} gives
\begin{equation}
\label{eq:grid-oscillation-bound}
 \frac{1}{r^2}\int_{B(x,r)}T_n\wedge\beta
 \leq C_R\left(
       \frac{\operatorname{osc}_{B(x,r)}\rho}{b-a}+\frac2n
             \right).
\end{equation}

For completeness, choose $\chi\in C_c^\infty(U)$ nonnegative and
equal to one on $B(x,R)$.  Extend $\chi\beta$ by zero to $S$.
The positive measures $\mu_n=T_n\wedge\chi\beta$ converge weakly
to $\mu=T\wedge\chi\beta$: current convergence gives convergence
on smooth test functions, and the uniform mass bound extends this
to continuous test functions.  The open-set inequality for weak
convergence of measures gives, for the open ball $B(x,r)$,
\[
 \int_{B(x,r)}T\wedge\beta
   =\mu(B(x,r))
   \leq\liminf_{n\to\infty}\mu_n(B(x,r)).
\]
Passing to the limit in \eqref{eq:grid-oscillation-bound} therefore
yields
\begin{equation}
\label{eq:limit-oscillation-bound}
 \frac{1}{r^2}\int_{B(x,r)}T\wedge\beta
 \leq \frac{C_R}{b-a}\operatorname{osc}_{B(x,r)}\rho.
\end{equation}
Continuity of $\rho$ makes the right side tend to zero as
$r\downarrow0$.  Thus $\nu(T,x)=0$.  At every point of $S\setminus X$
the current vanishes on a neighborhood, since its support is contained
in $K\Subset X$.  This proves \eqref{eq:averaged-current-properties}.
The same weakly convergent subsequence works for every fixed ball;
no diagonal choice depending on $x$ or $r$ is required.
\end{proof}

We record explicitly the consequence of the surface current theorem
that will be used.  In particular, a de Rham class is not assigned
to a nonclosed current without explanation.

\begin{lemma}[The surface current inequality]
\label{lem:surface-current-inequality}
Let $T$ be a positive $dd^c$-closed $(1,1)$-current on a compact
K\"ahler surface $S$, with $\nu(T,x)=0$ for all $x\in S$.
There is a well-defined class $\alpha_T\in H^{1,1}(S,\mathbb R)$
such that, for every smooth closed real $(1,1)$-form $\eta$,
\begin{equation}
\label{eq:current-class-pairing}
 \alpha_T\cdot[\eta]=\int_S T\wedge\eta.
\end{equation}
It satisfies
\begin{equation}
\label{eq:DNS-inequality}
 \alpha_T^2\geq0,
 \qquad \alpha_T^2>0\ \text{if }dT\ne0.
\end{equation}
If $dT=0$, then $\alpha_T$ is the usual cohomology class of $T$.
\end{lemma}

\begin{proof}
The decomposition in \cite[\S2.1, equation~(2.1), and Lemma~B.2]{DNS22}
writes
\[
 T=\Omega+\partial R+\bar\partial\overline R
          +i\partial\bar\partial u,
\]
where $\Omega$ is a smooth real closed $(1,1)$-form, $R$ is a
$(0,1)$-current, and $u$ is a real distribution.  Define
$\alpha_T=[\Omega]$.  Integration by parts gives
\eqref{eq:current-class-pairing}.  Nondegeneracy of the intersection
pairing on $H^{1,1}(S,\mathbb R)$ makes $[\Omega]$ independent of
the decomposition and identifies it with $[T]$ when $T$ is closed.

Since the positive-Lelong-number locus is empty,
\cite[Corollary~2.4 and its proof]{DNS22} gives
\eqref{eq:DNS-inequality}.  More specifically, its proof applies
\cite[Theorem~2.2(2)]{DNS22} to the pair $(T,T)$ and uses the
vanishing-energy criterion following equation~(B.4).  The hypothesis
there is that each current gives no mass to the positive-Lelong-number
locus of the other, which is automatic here.  These results concern
arbitrary positive $dd^c$-closed currents on compact K\"ahler surfaces;
they impose no foliation hypothesis.  We use only the qualitative
inequalities \eqref{eq:DNS-inequality}, which are independent of
the positive normalization factor in $dd^c$.
\end{proof}

\begin{proof}[Proof of Theorem~\ref{thm:divisor-regularization}]
Fix a K\"ahler form $\omega$ on $S$ and a continuous
plurisubharmonic exhaustion of $X$.  If $X$ is a modification
$\pi:X\to Y$ of a Stein space, a smooth plurisubharmonic
exhaustion of $Y$ pulls back to one on $X$, since $\pi$ is proper.
Smoothness on a complex space is understood in local embeddings;
Stein spaces admit smooth strictly plurisubharmonic exhaustions.

We may therefore suppose that $X$ is not a modification of a Stein
space.  Lemmas~\ref{lem:level-averaging} and
\ref{lem:surface-current-inequality} give a current $T$, supported
in a compact set $K\Subset X$, and a class
$\alpha=\alpha_T\in H^{1,1}(S,\mathbb R)$ with
\begin{equation}
\label{eq:alpha-basic}
 \alpha\cdot[\omega]=1,\qquad
 \alpha^2\geq0,\qquad
 \alpha^2>0\ \text{if }dT\ne0.
\end{equation}
In particular $\alpha\ne0$.

For each irreducible component $D_i$, choose a smooth Hermitian
metric $h_i$ on $\mathcal O_S(D_i)$, with curvature $\theta_i$,
and denote its canonical section by $s_i$.  Poincar\'e--Lelong gives
\[
 \ddc\log\|s_i\|_{h_i}^{2}=[D_i]-\theta_i.
\]
Choose a smooth cutoff equal to one on a neighborhood of $K$ and
supported away from $D_i$.  Multiplying $\log\|s_i\|_{h_i}^{2}$
by this cutoff and extending by zero gives a globally smooth
function $f_i$ with $\ddc f_i=-\theta_i$ near $K$.  Therefore
\begin{equation}
\label{eq:alpha-boundary-orthogonality}
 \alpha\cdot[D_i]
   =\int_S T\wedge\theta_i
   =-\langle T,\ddc f_i\rangle
   =0.
\end{equation}
This calculation uses only $dd^c$-closedness of $T$, not
$d$-closedness, and is legitimate by
\eqref{eq:current-class-pairing}.

We now treat all connected boundary components.  For such a component
$E$, let $I_E=(D_i\cdot D_j)_{D_i,D_j\subset E}$ and denote its
largest eigenvalue by $\lambda_E$.  If $\lambda_E<0$, the matrix
is negative definite, contrary to
Lemma~\ref{lem:surface-boundary-tools}\textup{(i)}.  If
$\lambda_E>0$, a real vector with positive quadratic value can
be approximated by a rational vector.  Clearing denominators gives
an integral divisor $B$ supported on $E$ with $B^2>0$.  (Its
coefficients need not be positive for this argument.)
Equation~\eqref{eq:alpha-boundary-orthogonality} gives
$\alpha\cdot[B]=0$.  The Hodge index theorem for the compact
K\"ahler surface says that the intersection form on
$H^{1,1}(S,\mathbb R)$ has signature $(1,h^{1,1}-1)$.
Consequently $[B]^\perp$ is negative definite, so
$\alpha^2<0$ because $\alpha\ne0$.  This contradicts
\eqref{eq:alpha-basic}.  We conclude that
\begin{equation}
\label{eq:boundary-top-eigenvalue}
 \lambda_E=0\qquad\text{for every connected component }E.
\end{equation}

The off-diagonal entries of $I_E$ are nonnegative, and their graph
is connected.  After adding a sufficiently large scalar multiple
of the identity, Perron--Frobenius applies to a nonnegative
irreducible matrix.  Thus the largest eigenvalue of $I_E$ is
simple and has an eigenvector all of whose entries are strictly
positive.  By \eqref{eq:boundary-top-eigenvalue}, its kernel is
exactly this one-dimensional eigenspace.  Since $I_E$ is integral,
its kernel has a rational generator.  Scaling the positive
generator therefore gives positive integers $a_i$ and a divisor
\[
 A_E=\sum_{D_i\subset E}a_iD_i,\qquad
 A_E\cdot D_i=0\quad(D_i\subset E),\qquad A_E^2=0.
\]
Set $A=\sum_E A_E$.  Distinct connected boundary components are
disjoint, so
\begin{equation}
\label{eq:full-support-null-divisor}
 \operatorname{supp}A=D,\qquad
 A\cdot D_i=0\ \text{for every }i,\qquad A^2=0.
\end{equation}
Every coefficient of $A$ is strictly positive.  None of these
matrix arguments requires the components to be smooth or the
divisor to have normal crossings.

The class $[A]$ is nonzero, since $A$ is nonzero and effective and
$\int_A\omega>0$.  Equations
\eqref{eq:alpha-boundary-orthogonality} and
\eqref{eq:full-support-null-divisor} give $\alpha\cdot[A]=0$.
For a nonzero isotropic vector in a real vector space of signature
$(1,h^{1,1}-1)$, its orthogonal complement is negative semidefinite
and has radical equal to the line it spans.  Hodge index therefore
gives $\alpha^2\leq0$, with equality only when $\alpha$ is a
multiple of $[A]$.  Combining this with \eqref{eq:alpha-basic}
proves
\begin{equation}
\label{eq:closed-current-null-class}
 dT=0,\qquad [T]=\alpha=c[A],\qquad
 c=\frac{1}{\int_A\omega}>0.
\end{equation}
Only now do we use $[T]$ as the de Rham class of a closed current.

The $dd^c$-lemma for currents on the compact K\"ahler manifold $S$
applied to \eqref{eq:closed-current-null-class} yields a real
distribution $u$ such that
\begin{equation}
\label{eq:global-log-potential}
 \ddc u=T-c[A].
\end{equation}
Choose a neighborhood $U$ of $D$ disjoint from $K$.  On a small
coordinate neighborhood in $U$, let $f_A$ be a local holomorphic
equation of the effective Cartier divisor $A$.  Then
\[
 \ddc\bigl(u+c\log|f_A|^2\bigr)=0
\]
as distributions on the entire coordinate neighborhood, including
its intersection with $D$.  A pluriharmonic distribution is smooth:
taking a trace gives a distributional Laplace equation, so elliptic
regularity applies.  Hence on that neighborhood
\begin{equation}
\label{eq:local-log-collar}
 u=-c\log|f_A|^2+h,
 \qquad h\text{ smooth and pluriharmonic}.
\end{equation}
These identities show that $u$ is a smooth pluriharmonic function
on $U\setminus D$ and tends to $+\infty$ as one approaches any
point of $D$.  The divergence is uniform near $D$ by compactness
and a finite coordinate covering.  At a normal crossing, for
example, the singular part is
$-c(a_i\log|z|^2+a_j\log|w|^2)$.  The same conclusion follows
from a local equation $f_A$ at an arbitrary singular point.

To see the asserted metric explicitly, take local frames $e_i$ of
$\mathcal O_S(A)$ with $s_A=f_i e_i$, and put
$b_i=r_i/c$, where $r_i=u+c\log|f_i|^2$ is the smooth
pluriharmonic remainder in \eqref{eq:local-log-collar}.  If
$e_j=g_{ij}e_i$, then $f_i=g_{ij}f_j$ and
\[
 b_j=b_i-\log|g_{ij}|^2.
\]
Thus $\|e_i\|^2=\exp(-b_i)$ defines a smooth Hermitian metric
on $\mathcal O_S(A)|_U$.  Its curvature is $\ddc b_i=0$, and
$-\log\|s_A\|^2=u/c$ on $U\setminus D$.

Finally choose $W\Subset U$ containing $D$, with compact boundary
disjoint from $D$.  Choose $M$ larger than the values of $u$ on a
neighborhood of $\partial W$.  Let $\chi:\mathbb R\to\mathbb R$
be smooth, convex, and nondecreasing, identically zero on
$(-\infty,M]$, and tending to $+\infty$ at $+\infty$.  Define
\[
 \Phi=
 \begin{cases}
   \chi(u),&\text{on }W\setminus D,\\
   0,&\text{on }X\setminus W.
 \end{cases}
\]
The function vanishes near $\partial W$, so it is globally smooth.
It is plurisubharmonic because $u$ is pluriharmonic on the collar
and $\chi$ is convex.  It diverges along all of $D$, by the
full-support condition in \eqref{eq:full-support-null-divisor}.
Compactness of $S$ then makes every sublevel of $\Phi$ compact
in $X$.  This completes the proof.
\end{proof}

\begin{corollary}[Smooth quasi-projective surfaces]
\label{cor:quasiprojective-regularization}
Let $X$ be a smooth complex quasi-projective surface.  If $X$
admits a continuous plurisubharmonic exhaustion, then it admits
a smooth plurisubharmonic exhaustion.
\end{corollary}

\begin{proof}
For connected compact $X$ the assertion is immediate.  Otherwise
take a projective closure $Y\subset\mathbb P^N$.  Embedded
resolution, chosen to preserve the regular locus of $Y$, gives a
smooth projective surface containing the given smooth open subset
$X$.  Principalize the ideal of its complement by blowups with
smooth centers supported outside $X$.  The full inverse-image
boundary is then an effective divisor with simple normal crossings,
and its complement is still isomorphic to $X$; see
\cite[Theorems~1.0.1--1.0.2]{Wlodarczyk05}.  Thus $X=S\setminus D$
for a smooth projective surface $S$ and a nonempty reduced divisor
$D$.  Apply Theorem~\ref{thm:divisor-regularization}.  The use of
the full inverse image includes any originally zero-dimensional
boundary.  For a disconnected $X$, apply the argument to its
finitely many connected components.
\end{proof}

\begin{remark}[Quasi-projective surfaces with only constant holomorphic functions]
\label{rem:quasiprojective-constant-functions}
Corollary~\ref{cor:quasiprojective-regularization} is not subsumed by
Proposition~\ref{prop:surface-nonconstant-function}: a noncompact
smooth quasi-projective surface can admit a smooth psh exhaustion
while satisfying $\OO(X)=\C$, where $\OO(X)$ denotes all analytic
holomorphic functions, not merely algebraic regular functions.
Indeed, let $B$ be a smooth projective curve of positive genus,
let $A\in\operatorname{Pic}^0(B)$ be non-torsion, and put
$X=\Tot(A^{-1})$.  A unitary flat metric on $A^{-1}$ gives a smooth
psh exhaustion by its squared fiber norm, as in
Proposition~\ref{cor:curve-total-space}.

For any holomorphic function on $X$, its degree-$n$ fiberwise Taylor
coefficient along the zero section is a section of $A^n$.
For $n>0$ these sections vanish, since a nontrivial degree-zero line
bundle has no nonzero holomorphic section.  The degree-zero
coefficient is constant on $B$, so $\OO(X)=\C$.
The projective completion
$S=\PP_{\mathrm{lines}}(\OO_B\oplus A^{-1})$ has a single smooth
infinity section $D=\PP_{\mathrm{lines}}(A^{-1})$ with
$D^2=0$ and $N_{D/S}\simeq A$.

These examples form large families.  For fixed $B$ of genus $g\ge1$,
a Poincar\'e line bundle gives an algebraic family of such total
spaces over the $g$-dimensional abelian variety
$\operatorname{Pic}^0(B)$.  The torsion subgroup is countable, so
the preceding conclusion holds for every parameter outside a
countable set, hence for a very general member of this family.
This is a statement within the indicated family, not a genericity
assertion about all quasi-projective surfaces.

Moreover, let $\mu:\widetilde X\to X$ be the blowup of any finite
set of $n$ distinct points of $X$.  Then $\widetilde X$ is again
smooth and quasi-projective, and
$\mu_*\OO_{\widetilde X}=\OO_X$ gives
$\OO(\widetilde X)=\C$.  The pullback of the squared-norm exhaustion
is smooth and psh, and remains proper because $\mu$ is proper.
Blowing up the same points in $S$ leaves the infinity section and its
normal bundle unchanged.  Since $X$ retracts onto $B$, the rational
Betti numbers of $\widetilde X$ are
\[
 b_1(\widetilde X)=2g,\qquad b_2(\widetilde X)=n+1.
\]
Thus both Betti numbers can be arbitrarily large within this class.

Thus the constant-function case occurs even for smoothly weakly
complete quasi-projective surfaces with a connected square-zero
boundary.  Quasi-projectivity supplies the divisor compactification
used in our argument; it does not supply the nonconstant holomorphic
function required by
Proposition~\ref{prop:surface-nonconstant-function}.
\end{remark}

\begin{corollary}[Continuous collars along square-zero curves]
\label{cor:square-zero-continuous-collar}
Let $C$ be a smooth connected curve with $C^2=0$ in a smooth compact
connected K\"ahler surface $S$.  Suppose a neighborhood $V$ of $C$ admits
a continuous plurisubharmonic function on $V\setminus C$ tending
to $+\infty$ along $C$.  Then $S\setminus C$ admits a smooth
plurisubharmonic exhaustion.  If $S\setminus C$ is not a
modification of a Stein space, $\mathcal O_S(C)$ is unitary flat
near $C$, and the divergent collar can be chosen pluriharmonic.
No genus or Ueda-type hypothesis is required.
\end{corollary}

\begin{proof}
Choose a smaller neighborhood of $C$ compactly contained in $V$.
The given function is bounded on a compact annulus at its outer
boundary.  Taking its maximum with a larger constant and extending
by that constant gives a continuous plurisubharmonic exhaustion
of $S\setminus C$.  Apply
Theorem~\ref{thm:divisor-regularization}.  In its second alternative,
the boundary consists of the single component $C$, so the positive
kernel vector may be chosen to be $1$, giving $A=C$.  A smooth
zero-curvature Hermitian metric is unitary flat.  Equivalently, on
simply connected coordinate neighborhoods, write the smooth
pluriharmonic remainder in \eqref{eq:local-log-collar} as
$h/c=2\operatorname{Re}H$.  Replacing a defining function $w$ by
$w\exp(-H)$ produces defining functions of equal modulus on
overlaps, whose ratios are locally constant and unitary.
\end{proof}

\begin{remark}[Scope and the weighted boundary]
\label{rem:regularization-scope}
For a reducible boundary the flat metric is on the weighted line
bundle $\mathcal O_S(A)$; flatness of
$\mathcal O_S(D)$ is not asserted.  Positivity of every coefficient
of $A$ is what makes the logarithmic function exhaust every end.
The argument proves existence of a new smooth exhaustion and does
not assert smooth approximation of the original continuous function.
It covers all smooth quasi-projective surfaces and, more generally,
divisor complements in compact K\"ahler surfaces.  It does not
provide a divisorial compactification of an arbitrary noncompact
K\"ahler surface, or treat an arbitrary analytic domain in a
projective surface whose boundary is not a divisor.
\end{remark}


\subsection{Nonprojective compactifiers}\label{sec:surface-nonprojective}

\begin{theorem}[Nonprojective boundary criterion]\label{thm:nonprojective}
Let $S$ be a smooth compact connected nonprojective surface and let
$D\ne0$ be effective. Then the following are equivalent:
\begin{enumerate}[label=\textup{(\roman*)},leftmargin=2.5em]
\item $S\setminus\supp D$ has a continuous psh exhaustion;
\item $S\setminus\supp D$ has a smooth psh exhaustion;
\item no connected component of $\supp D$ has negative definite
intersection matrix.
\end{enumerate}
No K\"ahler or SNC hypothesis is required. The Hopf branch of the
proof uses the elliptic-neighborhood theorem in the preprint
\cite[Theorem~1.6]{CFSTZ2026}.
\end{theorem}
\begin{proof}
Only \textup{(iii)}$\Rightarrow$\textup{(ii)} needs proof, by
Lemma~\ref{lem:surface-boundary-tools}\textup{(i)}. We separate the
algebraic-dimension cases, using the classification facts in
\cite[Section~2]{FN2005}.

If $a(S)=1$, the algebraic reduction is a holomorphic elliptic
fibration $f:S\to B$ with connected fibers, and every curve is
vertical. For a fiber $F=\sum m_iC_i$, its matrix $I$ satisfies
\begin{equation}\label{eq:surface-fiber-intersection-identity}
 x^{\mathsf T}Ix=-\sum_{i<j}(C_i\cdot C_j)m_im_j
       \left(\frac{x_i}{m_i}-\frac{x_j}{m_j}\right)^2.
\end{equation}
The kernel is the line spanned by $(m_i)$, and every proper principal
submatrix is negative definite. Hence the boundary contains every
fiber support it meets. Thus $X=f^{-1}(B\setminus P)$ for a finite
nonempty set $P$, and a smooth strictly subharmonic exhaustion of the
punctured compact curve pulls back to a smooth psh exhaustion of $X$.
This argument includes reducible and multiple fibers.

Suppose $a(S)=0$. For every line bundle $L$ on a nonprojective surface,
$c_1(L)^2\le0$ \cite[Lemma~2.1]{FN2005}. Every boundary matrix is
therefore negative semidefinite. Hypothesis \textup{(iii)} and the
same Perron--Frobenius argument give an effective divisor $A$ with
positive integral coefficients, full boundary support, and $A^2=0$.
Contract to a smooth minimal model $p:S\to S_0$. At each blowdown
$q:T\to T'$, write $A=q^*A'+bE$, where $A'=q_*A$ and $E^2=-1$.
Then
\[
 0=A^2=(A')^2-b^2,\qquad (A')^2\le0,
\]
so $b=0$ and $A'\ne0$. Consequently $A=p^*A_0$, $A_0^2=0$, and
\[
 \supp A=p^{-1}(\supp A_0).
\]
The restriction $X\to S_0\setminus\supp A_0$ is proper. It is enough
to construct a smooth exhaustion downstairs.

In the K\"ahler case, $S_0$ is a torus or a K3 surface of algebraic
dimension zero. A torus of algebraic dimension zero has no curves:
an irreducible curve would be elliptic (by adjunction and the absence
of rational curves), and its elliptic quotient would give positive
algebraic dimension. On a K3 surface, Riemann--Roch gives
$h^0(\OO(A_0))\ge2$ because $A_0^2=0$ and $-A_0$ is not effective.
Both alternatives contradict $a(S_0)=0$.

In the non-K\"ahler case, $S_0$ is of class~VII. If $b_2(S_0)=n>0$,
Enoki's Main Theorem \cite[p.~454]{Enoki} identifies a nonzero
effective square-zero divisor with a multiple of the distinguished
cycle on an Enoki surface. Its complement is an affine-line bundle
of degree $-n$ over an elliptic curve \cite[(0.4), p.~453]{Enoki}.
Complete the affine fibers to a projective $\PP^1$-bundle. Its infinity
section has square $n>0$, so
Lemma~\ref{lem:surface-boundary-tools}\textup{(iii),(ii)} gives a smooth
psh exhaustion of the affine-line complement. This uses Enoki's
specific square-zero-divisor theorem, not a classification of all
class~VII surfaces.

If $b_2(S_0)=0$, a minimal class~VII surface containing a curve is
elliptic or Hopf \cite[p.~453]{Enoki}; algebraic dimension zero leaves
Hopf. Each irreducible curve has arithmetic genus one, since
$C^2=K_{S_0}\cdot C=0$. It cannot have rational normalization: a map
from $\PP^1$ lifts to $\C^2\setminus\{0\}$ and is constant. Thus the
curves are smooth elliptic and mutually disjoint. For each of them,
finite Ueda type gives a smooth strictly psh divergent collar
\cite[Theorem~2.5]{CLPT2018}. Infinite type gives a unitary flat
neighborhood by \cite[Theorem~1.6]{CFSTZ2026}, and hence the smooth
pluriharmonic collar $-\log|w_i|^2$ in unitary defining functions.
Lemma~\ref{lem:surface-boundary-tools}\textup{(ii)} finishes the proof.
\end{proof}

In particular, for a nonprojective K\"ahler compactifier, a nonempty
boundary complement can be continuously weakly complete only when
$a(S)=1$ and the boundary is a union of whole fiber supports. The
$a(S)=0$ K\"ahler branch is excluded, not an unresolved positive case.

\subsection{Smooth rulings}\label{subsec:smooth-rulings}

\begin{theorem}[Boundary criterion for a smooth ruling]
\label{cor:ruled-all-degrees}
Let $f:S\to B$ be a smooth $\PP^1$-bundle over a compact connected
Riemann surface, and let $D\ne0$ be reduced SNC. Then continuous psh
exhaustion, smooth psh exhaustion, and the absence of a negative
definite connected boundary component are equivalent for $S\setminus D$.
There is no restriction on the degree $D\cdot F$.
\end{theorem}
\begin{proof}
Necessity follows from Lemma~\ref{lem:surface-boundary-tools}. For
sufficiency, positive largest eigenvalue gives a smooth collar. A
holomorphic $\PP^1$-bundle over a compact curve is a projectivized
rank-two bundle; in particular, it is projective and has a section.
For a
connected component $E=\bigcup E_i$ with zero largest eigenvalue, take
$A=\sum a_iE_i$ as in \eqref{eq:surface-null-boundary}. We construct a
smooth exhaustion of $S\setminus E$, whose restriction supplies the
required collar at $E$.

If every component is vertical, $E$ is one whole fiber, since the
fibers of a smooth ruling are irreducible. Pull back an exhaustion
of $B$ minus its image. Otherwise $A\cdot F>0$, so no fiber can be a
component of $E$. Choose a section $C_0$, with $C_0^2=-e$; the
numerical divisor space has basis $C_0,F$. Writing
$A\equiv aC_0+bF$, the equality $A^2=0$ gives $b=ae/2$.
Writing $E_i\equiv d_iC_0+b_iF$, the identities $A\cdot E_i=0$ give
\[
 E_i\equiv d_i\bigl(C_0+(e/2)F\bigr),\qquad
 E_i^2=E_i\cdot E_j=0\quad(i\ne j).
\]
Distinct components are therefore disjoint. Connectedness and the
SNC hypothesis make $E=C$ a single smooth square-zero curve. For
$d=C\cdot F$, adjunction and
$K_S\equiv-2C_0+(2g(B)-2-e)F$ give
\[
 2g(C)-2=d(2g(B)-2).
\]
Riemann--Hurwitz shows that $C\to B$ is finite \'etale.

If $d=1$, the complement is an affine-line bundle. Its linear bundle
is dual to $N_{C/S}$ and therefore has degree zero. We recall the smooth exhaustion construction of
Koike--Ueda \cite[Proposition~2.1]{KoikeUeda2025}. In unitary flat
frames its transitions are $w_i=a_{ij}w_j+b_{ij}(z)$, $|a_{ij}|=1$.
Choose a smooth section and subtract a global smooth section of the linear
bundle so that its $\bar\partial$ is harmonic. On coordinate discs,
choose antiholomorphic primitives $g_i$ of this harmonic form.
Changing the affine coordinates by holomorphic translations then gives
\[
 w_i=a_{ij}w_j+c_{ij},\qquad
 g_i=a_{ij}g_j+c_{ij},
\]
where $c_{ij}$ is constant: the differences of the $g_i$ are both
holomorphic and antiholomorphic. The functions
$\rho_i=|w_i-g_i|^2$ agree and have Levi matrix
\[
 \begin{pmatrix}|(g_i)_{\bar z}|^2&0\\0&1\end{pmatrix}.
\]
They define a smooth psh exhaustion; its sublevels are bounded closed
disc bundles around a smooth section over compact $B$.

For $d\ge2$, a finite \'etale Galois base change splits the boundary
into $d$ disjoint sections. Each has square zero, since the pullback
of $C$ is their disjoint sum and the projection formula applies.
If $d=2$, putting these sections at $0$ and $\infty$ identifies the
complement with $L^\times$ for a degree-zero line bundle. A unitary
flat metric makes $(\log\|v\|^2)^2$ a smooth psh exhaustion.
If $d\ge3$, three labeled sections trivialize the $\PP^1$-bundle;
the others are constant, as maps from compact $B'$ to $\PP^1$
avoiding $0,1,\infty$. The complement is
$B'\times(\PP^1\setminus Q)$ for a finite set $Q$, and a smooth
subharmonic exhaustion of the punctured projective line pulls back.
Average a nonnegative exhaustion over the finite deck group. The
average remains proper, being at least the original exhaustion divided
by the group order, and descends smoothly to $S\setminus C$.
Finally glue the collars at all connected boundary components using
Lemma~\ref{lem:surface-boundary-tools}\textup{(ii)}.
\end{proof}

This numerical existence criterion uses the actual smooth ruling,
not merely birational ruledness. Regularization for arbitrary smooth
projective blowups and singular rational fibrations is nevertheless
covered by Theorem~\ref{thm:divisor-regularization}; the numerical
criterion alone is not asserted for those compactifiers. For $S=\PP^2$,
every nonzero effective divisor is ample, so its complement has a
smooth psh exhaustion directly from a positive metric on $\OO_S(D)$.

\subsection{Direct geometric constructions in the projective case}
\label{subsec:projective-cases}

The conditional regularization conclusions in the next proposition
are now special cases of Theorem~\ref{thm:divisor-regularization}.
We retain their proofs because they give direct geometric
constructions from disjoint boundary divisors, moving divisors,
surface abundance, and fibrations.


\begin{proposition}[Projective affirmative cases]
\label{prop:surface-projective-cases}
Let $S$ be smooth projective, let $D\ne0$ be reduced SNC, and assume
$X=S\setminus D$ has a continuous psh exhaustion. Then $X$ has a
smooth psh exhaustion in each of the following cases:
\begin{enumerate}[label=\textup{(\roman*)},leftmargin=2.5em]
\item $D$ is disconnected;
\item on every connected boundary component with zero largest
intersection eigenvalue, a positive integral kernel divisor $A$
satisfies $\kappa(S,A)=1$ or $K_S\cdot A<0$;
\item $K_S$ is torsion in $\operatorname{Pic}(S)$;
\item every component of $D$ is vertical for a morphism
$f:S\to B$ with connected fibers onto a smooth projective curve.
\end{enumerate}
\end{proposition}
\begin{proof}
For \textup{(i)}, write the connected boundary components as
$E_1,\ldots,E_r$, $r\ge2$. The matrix argument provides effective
$A_k$ with support $E_k$ and $A_k^2\ge0$, with strict inequality in
the positive-eigenvalue case. Since $A_k\cdot A_l=0$ for $k\ne l$,
the Hodge index theorem excludes a positive square and makes the
nonzero square-zero nef classes proportional. The proportionality
factors are positive rational numbers, by intersection with an
integral ample class. Choose positive integers $p_k,q_k$ such that
$p_kA_k-q_kA_1$ is numerically trivial, for $k\ge2$. Its real first
Chern class is zero by Hodge index, and the $\partial\bar\partial$-lemma
gives a unitary flat metric $h_k$ on
$\OO_S(p_kA_k-q_kA_1)$. For its canonical meromorphic section
$\sigma_k$, the functions
\[
 \phi_k=\log\|\sigma_k\|_{h_k}^2,\qquad
 \rho=\sum_{k=2}^r\phi_k^2
\]
are smooth on $X$; the $\phi_k$ are pluriharmonic, so $\rho$ is psh.
It tends to infinity at every boundary component, proving properness.

For \textup{(ii)}, a nef divisor of Iitaka dimension one on a surface
is semiample \cite[Lemma~7.3]{Fujino2021}. If $K_S\cdot A<0$,
Riemann--Roch gives
\[
 \chi(\OO_S(nA))=\chi(\OO_S)-\tfrac n2K_S\cdot A\longrightarrow+\infty.
\]
Serre duality and intersection with an ample class give
$H^2(S,\OO_S(nA))=0$ for large $n$. Thus $\kappa(S,A)\ge1$, and
nefness with $A^2=0$ gives $\kappa(S,A)=1$. In either case $A$ is
semiample. A base-point-free multiple defines, after Stein
factorization, a morphism $f:S\to B$ to a smooth curve with connected
fibers and $\OO_S(mA)=f^*L$. The canonical section descends because
$f_*\OO_S=\OO_B$, so $mA=f^*G$ for a nonzero effective divisor $G$
on $B$. Hence
\[
 S\setminus\supp A=f^{-1}(B\setminus\supp G),
\]
with proper restricted map. Pulling back an exhaustion of the
punctured curve gives a collar at $\supp A$. Combine these with the
positive-eigenvalue collars using
Lemma~\ref{lem:surface-boundary-tools}\textup{(ii)}.

For \textup{(iii)}, consider such a square-zero nef $A$. Choose a
small positive rational $\varepsilon$ so that every coefficient of
$\varepsilon A$ is at most one. The log divisor
$K_S+\varepsilon A$ is nef. Log abundance for surfaces
\cite[Theorem~6.1]{Fujino2012} makes it semiample. Killing the torsion
of $K_S$ and clearing denominators shows that $A$ itself is semiample.
Apply the preceding fibration argument. This is a direct consequence
of log abundance on the given pair; no change of the open surface by
a log minimal model program is involved.

For \textup{(iv)}, identity
\eqref{eq:surface-fiber-intersection-identity} holds for the
multiplicity vector of any fiber. Every proper subset of its
components has negative definite intersection matrix. Therefore the
boundary obstruction forces $D$ to be a union of whole fiber supports.
The restricted map $X\to B\setminus P$ is proper for a finite nonempty
set $P$; pull back an exhaustion of the punctured base.
\end{proof}

In \textup{(i)} no SNC assumption is necessary, because boundary
resolution preserves the complement and its number of boundary
components. Part \textup{(iii)} applies, in particular, to the minimal
projective surfaces of Kodaira dimension zero: K3, Enriques, abelian,
and bielliptic surfaces \cite[Chapter~VI]{BHPV2004}. This is a standard
consequence of the cited abundance theorem, not a new classification
result. Part \textup{(iv)} includes the vertical-boundary case of an
elliptic fibration, while arbitrary horizontal boundaries under the continuous-exhaustion
hypothesis are covered by Theorem~\ref{thm:divisor-regularization}.

The weighted square-zero configurations that arise from the
boundary matrix are all covered by Theorem~\ref{thm:divisor-regularization}
under the continuous-exhaustion hypothesis. They therefore do not
constitute a remaining projective counterexample class.

\subsection{A concise smooth-boundary criterion}
\label{subsec:smooth-boundary}

The following sufficient conditions remain useful because no continuous
exhaustion is assumed. If a continuous collar is already given, the
unrestricted square-zero regularization statement is
Corollary~\ref{cor:square-zero-continuous-collar}.

For a smooth square-zero curve $C$, its normal bundle $N$ is unitary
flat. The Ueda type measures the orders to which local defining
functions can be made to have unitary linear transition functions.
Infinite type is formal; a unitary flat neighborhood requires
\emph{convergent} defining functions $w_i=t_{ij}w_j$, $|t_{ij}|=1$.
These two conditions must not be conflated \cite{Ueda1983,CLPT2018}.

\begin{proposition}[Sufficient conditions, not a classification]
\label{prop:surface-ueda-sufficient}
Let $S$ be a compact K\"ahler surface and let $C\subset S$ be smooth,
connected, and square-zero. Write $\rho:H_1(C,\mathbb Z)\to\Uone$
for the unitary character of $N_{C/S}$ and
$K=\ker[H_1(C,\mathbb Z)\to H_1(S,\mathbb Z)]$.
Then $S\setminus C$ has a smooth psh exhaustion if any of the following
holds: finite Ueda type; a unitary flat neighborhood; torsion or
Diophantine normal bundle; $\rho(K)$ finite; $q(S)\ge g(C)$; or
$g(C)=1$. No continuous exhaustion is assumed. Only the last case
uses \cite[Theorem~1.6]{CFSTZ2026}.
\end{proposition}
\begin{proof}
Finite type gives a smooth strictly psh collar
\cite[Theorem~2.5]{CLPT2018}. A flat neighborhood gives the
pluriharmonic collar $-\log|w_i|^2$. In infinite type, torsion or
Diophantine $N$ gives such a neighborhood by Ueda's convergence
theorem \cite[Theorem~3]{Ueda1983}; finite type has already been
handled. Here Diophantine means
$\operatorname{dist}(N^m,\OO_C)\ge c m^{-\tau}$ for some $c,\tau>0$
and all $m\ge1$.

If $\rho(K)$ is finite, a power $\rho^k$ is trivial on $K$, hence
defines a character of the image of $H_1(C,\mathbb Z)$ in
$H_1(S,\mathbb Z)$. It extends to the whole group because $\Uone$
is divisible. In the remaining non-torsion, infinite-type case,
Neeman's extension theorem in \cite[Theorem~4.4]{CLPT2018} gives an
effective divisor $E$ disjoint from $C$ and cohomologous to $mC$.
A unitary flat metric on $\OO_S(mC-E)$ then gives the pluriharmonic
collar $-m^{-1}\log\|\sigma\|^2$ at $C$.

For $q(S)\ge g(C)$, consider
$r:H^1(S,\OO_S)\to H^1(C,\OO_C)$. If $r$ is not injective,
\cite[Proposition~2.7]{CLPT2018} makes a multiple of $C$ a fiber of a
fibration, and the punctured-base exhaustion applies. Otherwise the
dimension inequality makes $r$ an isomorphism. Hodge decomposition
and duality make $H_1(C,\mathbb Q)\to H_1(S,\mathbb Q)$ an
isomorphism. Its integral kernel is torsion and lies in the free
abelian group $H_1(C,\mathbb Z)$, so $K=0$. Apply the character case.
Finally, when $g(C)=1$, infinite type implies a unitary flat
neighborhood by \cite[Theorem~1.6]{CFSTZ2026}. In every case finish
with Lemma~\ref{lem:surface-boundary-tools}\textup{(ii)}.
\end{proof}

The finite-type and flat-neighborhood arguments also apply without a
K\"ahler hypothesis, as used in the Hopf branch above. Neither
semipositivity of $\OO_S(C)$ nor holomorphic convexity of $S\setminus C$
is asserted as a necessary condition: the proposition constructs
collars, not a general equivalence with either of those properties.

\subsection{Classification of compactifiers and the scope of the theorem}
\label{subsec:surface-classification}
In Table~\ref{tab:surface-status}, $S$ is a smooth compact connected
surface, $D\ne0$ is a divisor, and $X=S\setminus\supp D$ is assumed
to have a continuous psh exhaustion. The familiar minimal types are
listed to organize the compactifiers; the projective conclusion
holds for every smooth projective model, including nonminimal ones.
``Affirmative'' means that a smooth psh exhaustion exists;
``excluded'' means that no such nonempty divisor complement satisfies
the continuous-exhaustion hypothesis. The table uses the
Enriques--Kodaira classification \cite[Chapter~VI]{BHPV2004}, and the
individual exhaustion conclusions follow from the theorems above.

All projective Kodaira dimensions are covered by
Theorem~\ref{thm:divisor-regularization}. In particular, rational
blowups, horizontal boundaries on properly elliptic surfaces, and
connected square-zero boundaries on surfaces of general type require
no additional hypothesis. The proof works on the given compactifier;
it does not contract curves meeting the open surface. The
nonprojective existence criterion likewise includes nonminimal
compactifiers through its full-pullback argument.

The high-genus, infinite-Ueda-type collar problem is covered by
Corollary~\ref{cor:square-zero-continuous-collar}. When the complement
is not a modification of Stein, the current argument gives a smooth
flat metric near the boundary; otherwise the Stein modification gives
the smooth exhaustion. Neither infinite Ueda type alone nor the
numerical square-zero condition is asserted to imply existence of a
continuous exhaustion.

The unrestricted continuous-to-$C^2$ problem for complex surfaces,
including arbitrary noncompact K\"ahler surfaces without a divisor
compactification, is not resolved here. Nor does the theorem treat
every analytic domain in a projective surface: a real or otherwise
nondivisorial boundary is outside its hypotheses. The new
compactifiable K\"ahler regularization theorem is independent of the
2026 elliptic-neighborhood input used in the separate existence
criteria, and independent of the threefold constructions.

\begin{table}[!htbp]
\caption{Divisor complements by compactifier type; a continuous psh exhaustion is assumed. Projective conclusions include nonminimal models.}
\label{tab:surface-status}
\small
\renewcommand{\arraystretch}{1.16}
\begin{tabularx}{\textwidth}{@{}>{\raggedright\arraybackslash}p{0.29\textwidth}>{\raggedright\arraybackslash}X@{}}
\toprule
Compactifier $S$ & Exhaustion conclusion and exact scope\\
\midrule
Projective, $\kappa(S)=-\infty$: rational or ruled
& \textbf{Affirmative.} Every divisor complement satisfying the hypothesis, including arbitrary smooth projective blowups; Theorem~\ref{thm:divisor-regularization}.\\
\addlinespace
Projective, $\kappa(S)=0$: K3, Enriques, abelian, bielliptic models
& \textbf{Affirmative.} All boundaries and nonminimal models; Theorem~\ref{thm:divisor-regularization}. The torsion-canonical construction remains a direct alternative on the corresponding minimal models.\\
\addlinespace
Projective, $\kappa(S)=1$: properly elliptic
& \textbf{Affirmative.} Horizontal as well as vertical boundaries; Theorem~\ref{thm:divisor-regularization}.\\
\addlinespace
Projective, $\kappa(S)=2$: general type
& \textbf{Affirmative.} All boundaries, including connected weighted square-zero divisors; Theorem~\ref{thm:divisor-regularization}.\\
\midrule
Nonprojective, $a(S)=1$
& \textbf{Affirmative.} The boundary is a union of whole elliptic fiber supports; Theorem~\ref{thm:nonprojective}. This includes the relevant Kodaira and elliptic surfaces.\\
\addlinespace
Nonprojective K\"ahler, $a(S)=0$: torus or K3 models
& \textbf{Excluded.} No nonempty divisor complement satisfies the hypothesis; see the proof of Theorem~\ref{thm:nonprojective}.\\
\addlinespace
Nonprojective, $a(S)=0$, class~VII, $b_2=0$ on the minimal model
& \textbf{Affirmative when a boundary exists.} The minimal model is Hopf; the elliptic-collar argument uses \cite{CFSTZ2026}. Blowups of curve-free minimal models have only exceptional boundaries, excluded by negative definiteness.\\
\addlinespace
Nonprojective, $a(S)=0$, class~VII, $b_2>0$ on the minimal model
& \textbf{Affirmative for every admissible boundary.} The matrix condition forces the Enoki square-zero-divisor case; other boundaries are excluded by the negative definite obstruction.\\
\bottomrule
\end{tabularx}
\end{table}
\FloatBarrier
\clearpage

\section{Holomorphic convexity and regularization for curve fiberings}
\label{sec:curve-fibers}
This section gives a regularization criterion independent of the
compactification hypotheses in Section~\ref{sec:surfaces}.
Complex spaces are reduced, Hausdorff, and have countable topology.
Pseudoconvexity means the existence of a continuous real-valued
psh exhaustion, as in \cite[p.~455]{VijituForum2022}.
All fiber dimensions and irreducible components refer to the reduced
supports of the fibers.

\begin{theorem}\label{thm:curve-fiber-convexity}
Let \(X\) be a complex space of pure dimension \(n\ge2\), admitting a
continuous real-valued psh exhaustion \(\rho\).
Suppose \(F:X\to\C^{n-1}\) is holomorphic and every nonempty fiber
has dimension at most one.  Then \(X\) is holomorphically convex
and admits a smooth psh exhaustion.
No condition on \(\operatorname{Sing}X\) is required.
\end{theorem}

Neither normality nor Cohen--Macaulayness of \(X\) is assumed.
The map \(F\) may be nonproper and its fibers may be disconnected,
reducible, singular, or multiple.  The proof concerns arbitrary
singularities of the total space as well as singular fibers.
Throughout the proof, \(X,\rho,F\) satisfy the hypotheses of the theorem.

\subsection{Cycle-theoretic and convexity inputs}
We state the exact forms needed.
\begin{enumerate}[label=\textup{(C\arabic*)},ref=C\arabic*,leftmargin=*]
\item\label{curve:ext:gf}
\emph{Geometric flatness over a normal base.}
An equidimensional holomorphic map of reduced pure-dimensional complex
spaces to a normal complex space has an analytic family of effective
weighted fiber cycles, with support equal to the corresponding fiber.
Cycles may be noncompact and locally finite; the map need not be proper.
See \cite[Chapter~I, \S1(b), p.~816]{BarletMagnusson1998}, with the definitions on
pp.~814--815. See also \cite[Definition~1.1.1 and following
Remark~(ii); author preprint pp.~3--4]{BarletMagnusson2019}, referring to
\cite[Chapter~IV, Theorem~3.4.1]{BarletMagnusson2014}.
This is a statement about cycles, not flatness of the local rings.

\item\label{curve:ext:integrals}
\emph{Continuity of integration.}
For a continuous family \(Z_t\) of analytic \(d\)-cycles on a complex
space and a continuous compactly supported \((d,d)\)-form \(\omega\)
in the local-embedding sense, \(t\mapsto\int_{Z_t}\omega\)
is continuous.
See \cite[Chapter~IV, Proposition~2.3.1]{BarletMagnusson2014};
this exact location is also cited in
\cite[proof of Theorem~4.0.7, Step~4, author preprint p.~21]{BarletAlpha2017}.
Compact support on the ambient space supplies the required proper
support on the graph of the family.

\item\label{curve:ext:bishop}
\emph{Bishop compactness on a complex space.}
For any complex analytic space \(W\), compact effective \(d\)-cycles
whose supports lie in a fixed compact set and whose Hermitian volumes
are uniformly bounded form a relatively compact subset of the
compact-cycle space.
See \cite[Theorem~1, p.~26]{Barlet1999}.
The topology on pp.~25--26 has the following properties:
under convergence, containment of the limiting support in an open
set implies eventual containment of the approximating supports there;
local integration currents also converge.
The cycles need not meet \(W_{\rm reg}\).
\end{enumerate}

The following results are used in their stated complex-space forms
in \cite{VijituForum2022}.
\begin{enumerate}[label=\textup{(V\arabic*)},ref=V\arabic*,leftmargin=*]
\item\label{curve:ext:sublevels}
\emph{Psh sublevel criterion} (Proposition~2.2, p.~456).
For a continuous real-valued psh function \(u\) on a complex space \(W\),
the space \(W\) is holomorphically convex if and only if
\(\{u<t\}\) is holomorphically convex for every real \(t\).
\item\label{curve:ext:localglobal}
\emph{Convexity over a Stein base} (Proposition~2.4, p.~457).
A pseudoconvex space holomorphically convex over a Stein space is
holomorphically convex. Here convexity over the base means that each
base point has an open neighborhood with holomorphically convex
inverse image.
\item\label{curve:ext:collar}
\emph{Boundary-collar criterion} (Corollary~2.5, p.~457).
Let \(\pi:M\to S\) be holomorphic, with \(S\) Stein, and let
\(D\subset M\) be pseudoconvex and open.
If \(\overline D^{\,M}\to S\) is proper and an open neighborhood
\(U\) of \(\partial_M D\) has Stein reduced fibers under \(\pi\),
then \(D\) is holomorphically convex.
\end{enumerate}

\begin{remark}[Exhaustion clarification in \ref{curve:ext:collar}]
The printed proof of Corollary~2.5 calls a smooth auxiliary function
constructed across \(\partial_M D\) a relative exhaustion.
One must add the psh exhaustion of \(D\).
After shrinking to a relatively compact base neighborhood \(V\),
with closure in a larger such neighborhood, that construction gives
a smooth \(h\), defined near
\(\overline D\cap\pi^{-1}(\overline V)\), strictly psh near its
relative boundary. This set is compact, so \(h\) is bounded there.
For a continuous psh exhaustion \(\eta_D\) of \(D\), the function
\(\eta_D+h\) on \(D\cap\pi^{-1}(V)\) is a vertical exhaustion.
Outside the boundary collar, the relevant part of
\(\overline D\cap\pi^{-1}(\overline V)\) is a compact subset of \(D\);
the sum is bounded above there. Above this bound it is strictly psh.
This satisfies the definition of a \(1\)-convex map on pp.~456--457.
Theorem~2.3 and Proposition~2.4 of that paper complete the argument.
No smoothness of \(M\) or its fibers is required for this completion.
\end{remark}

We also use that a reduced complex curve without compact irreducible
components is Stein; see \cite[\S3, p.~1663]{Vajaitu2022}.
The empty space is included as Stein.

\subsection{Fiber cycles and stability}
\begin{lemma}[Uniform area of compact fiber components]\label{curve:lem:area}
Fix a Hermitian metric \(h\) on \(X\), defined using local embeddings.
For every compact \(L\subset X\) and \(z_0\in\C^{n-1}\), there are a
neighborhood \(B\) of \(z_0\) and \(A>0\) such that \(\int_C h\leq A\)
for every reduced compact irreducible fiber component \(C\subset L\)
over \(B\).
\end{lemma}
\begin{proof}
On each local irreducible branch of \(X\), a fiber is cut out by
\(n-1\) holomorphic functions. The analytic dimension theorem
(the generalized principal ideal theorem for local analytic rings)
gives dimension at least \(n-(n-1)=1\) for every component of this
zero set. The upper bound therefore makes every nonempty reduced
fiber pure one-dimensional. This also holds for the restriction to
each irreducible source component, as required by \ref{curve:ext:gf}.

Since the base is normal, \ref{curve:ext:gf} supplies analytic weighted
fiber cycles \(\mathcal F_z\) with
\(|\mathcal F_z|=F^{-1}(z)\). Every component has a positive
integer coefficient.
Choose a nonnegative smooth compactly supported function \(\chi\)
equal to one near \(L\).
By \ref{curve:ext:integrals}, the function
\[
 z\longmapsto\int_{\mathcal F_z}\chi h
\]
is continuous and hence bounded by some \(A\) near \(z_0\).
Positivity gives
\[
 \int_C h=\int_{[C]}\chi h
 \leq\int_{\mathcal F_z}\chi h\leq A.
\]
There is no need to identify these geometric multiplicities with
scheme-theoretic fiber multiplicities.
\end{proof}

\begin{lemma}[Stability of noncompact components]\label{curve:lem:stability}
Let \(K\subset X\) be compact. If every irreducible component of
\(F^{-1}(z_0)\) meeting \(K\) is noncompact, then the same holds
for every fiber over some neighborhood of \(z_0\).
\end{lemma}
\begin{proof}
The assertion is immediate if \(K\) is empty. If the central fiber
misses \(K\), a ball avoiding the compact set \(F(K)\) suffices.
Otherwise suppose the assertion fails. There are \(z_j\to z_0\),
compact irreducible components \(C_j\subset F^{-1}(z_j)\), and
\(x_j\in C_j\cap K\). Pass to a subsequence with \(x_j\to x\in K\).
The function \(\rho\) is constant on each \(C_j\):
its pullback to the compact connected Riemann surface normalizing
\(C_j\) is subharmonic, hence constant. Thus
\[
 C_j\subset L:=\{\rho\leq\max_K\rho\},
\]
a fixed compact set.
Lemma~\ref{curve:lem:area} bounds their areas, so \ref{curve:ext:bishop}
gives a subsequence \([C_j]\) converging to a compact effective
one-cycle \(T\).

The limit is nonzero: otherwise upper support-continuity applied
to the empty open set forces \(C_j\) eventually to be empty.
Moreover \(x\in|T|\). If not, choose an open neighborhood of the
compact \(|T|\) whose closure misses \(x\). It eventually contains
every \(C_j\), contradicting \(x_j\to x\).

Also \(|T|\subset F^{-1}(z_0)\). A point outside the central fiber
has a neighborhood on which \(F\) avoids a neighborhood of \(z_0\).
This neighborhood misses every sufficiently large \(C_j\);
local convergence of currents shows that \(T\) has no support there.
An irreducible component of \(|T|\) through \(x\) is now a compact
closed irreducible analytic curve in the pure one-dimensional
central fiber. A closed irreducible analytic subset of the same
dimension as an irreducible curve containing it equals that curve.
It is therefore a full irreducible component of the central fiber,
contradicting the hypothesis.
\end{proof}

\begin{remark}
This replaces the smooth-source Corollary~4.2 of
\cite[p.~462]{VijituForum2022}.
The curves \(C_j\), or their limit, may lie entirely in
\(\operatorname{Sing}X\).
Multiple fibers are allowed because effective weighted fiber cycles
dominate their reduced component cycles.
\end{remark}

\subsection{Convex neighborhoods and globalization}
\begin{lemma}[Convex sublevel neighborhoods]\label{curve:lem:neighborhood}
For every \(\lambda\in\mathbb R\) and \(z_0\in\C^{n-1}\), there
are \(c>\lambda\) and a ball \(B=B_r(z_0)\) such that
\(D=F^{-1}(B)\cap\{\rho<c\}\) is holomorphically convex.
Consequently \(F^{-1}(B)\cap\{\rho<\lambda\}\) is holomorphically convex.
\end{lemma}
\begin{proof}
Let \(S\) be the set of constant \(\rho\)-values on the compact
irreducible components of \(F^{-1}(z_0)\).
For each real \(R\), only finitely many such components have value
at most \(R\): they meet the compact set \(\{\rho\leq R\}\),
while the irreducible components of a fixed analytic set are
locally finite. Thus \(S\) is locally finite.
Choose
\[
 \lambda<a<b,\qquad [a,b]\cap S=\varnothing,\qquad
 K=\{a\leq\rho\leq b\}.
\]
Every central irreducible component meeting \(K\) is noncompact.
Lemma~\ref{curve:lem:stability} gives a ball \(B=B_r(z_0)\) with the
same property for all fibers over \(B\).
Choose \(c\in(a,b)\) and put
\[
 M=F^{-1}(B),\qquad D=M\cap\{\rho<c\},\qquad
 U=M\cap\{a<\rho<b\}.
\]
Then \(\partial_M D\subset U\).

No reduced curve \(U\cap F^{-1}(z)\) has a compact irreducible
component \(A\). Such an \(A\), being compact and analytic in \(U\),
would be a closed analytic curve in \(X\). It would therefore be
a full irreducible component of the pure one-dimensional full fiber,
meeting \(K\), a contradiction.
Thus the Stein-curve criterion makes all reduced \(U\)-fibers Stein.

A continuous psh exhaustion of \(D\) is
\[
 \eta(x)=\max\left\{\frac1{c-\rho(x)},
                    \frac1{r^2-\|F(x)-z_0\|^2}\right\}.
\]
Both terms are psh by convex increasing composition.
For \(t>0\), its closed \(t\)-sublevel lies in
\[
 \{\rho\leq c-1/t\}\cap
 F^{-1}\{\|z-z_0\|^2\leq r^2-1/t\},
\]
is closed there, and stays inside \(D\); hence it is compact.
For \(t\leq0\) it is empty.
Furthermore, \(\overline D^{\,M}\to B\) is proper:
over a compact \(T\subset B\) its inverse image is a closed
subset of the compact set \(\{\rho\leq c\}\cap F^{-1}(T)\).
The boundary-collar criterion \ref{curve:ext:collar} proves that \(D\)
is holomorphically convex.
The final assertion follows from \ref{curve:ext:sublevels} for \(\rho|_D\).
\end{proof}

\begin{proof}[Proof of Theorem~\ref{thm:curve-fiber-convexity}]
For \(X(\lambda)=\{\rho<\lambda\}\),
Lemma~\ref{curve:lem:neighborhood} says that \(F|_{X(\lambda)}\) is a
holomorphically convex map to \(\C^{n-1}\).
The space \(X(\lambda)\) is pseudoconvex with exhaustion
\(1/(\lambda-\rho)\).
Criterion~\ref{curve:ext:localglobal} therefore makes it holomorphically
convex. This holds for every \(\lambda\), so
\ref{curve:ext:sublevels} proves that \(X\) is holomorphically convex.

Let \(q:X\to Z\) be its proper Remmert reduction, with \(Z\) Stein.
We use the standard implication from holomorphic convexity to smooth
weak completeness, recalled in \cite[p.~1659]{Vajaitu2022}.
Choose a smooth strictly psh exhaustion \(u\) of the Stein complex
space \(Z\).  As in Proposition~\ref{prop:surface-nonconstant-function},
the pullback \(u\circ q\) is smooth in the local-embedding sense
and psh.  It is exhaustive because
\[
 \{u\circ q\le t\}=q^{-1}(\{u\le t\})
\]
is compact for every real \(t\).  This proves the smooth-exhaustion
assertion and also Theorem~\ref{thm:curve-fiber-introduction}.
\end{proof}

\begin{corollary}[Simultaneous contraction]\label{cor:curve-contraction}
There are a Stein space \(Z\) and a proper surjective holomorphic
map \(q:X\to Z\), with connected fibers and \(q_*\OO_X=\OO_Z\),
whose nontrivial fibers are the maximal connected unions of
compact irreducible fiber curves of \(F\).
The map \(F\) factors holomorphically through \(q\).
\end{corollary}
\begin{proof}
Take the Remmert reduction of \(X\), using the reduction theory
recalled in \cite[proof of Proposition~2.4, p.~457]{VijituForum2022}.
The sheaf property makes the components of \(F\) descend, so each
\(q\)-fiber lies in an \(F\)-fiber and has dimension at most one.
A positive-dimensional compact connected analytic set of dimension
at most one is a finite connected union of compact irreducible curves.
Each is a full component of the corresponding pure one-dimensional
\(F\)-fiber.
Conversely, every compact irreducible curve is mapped to a point
in the Stein space \(Z\); hence so is any connected chain of them.
Such a chain has constant \(\rho\) and lies in a compact sublevel.
Local finiteness of the components of its \(F\)-fiber therefore
makes each maximal connected union finite.
This identifies the fibers of \(q\).
\end{proof}

\begin{remark}
For arbitrary reduced \(X\), normality of \(Z\) is not asserted;
it holds if \(X\) is normal.
Global holomorphic convexity supplies the contraction, without
requiring an independent gluing of locally defined quotients.
\end{remark}

\begin{remark}[Scope of the regularization conclusion]
The theorem constructs a smooth psh exhaustion; it does not
approximate the prescribed continuous exhaustion or assert that
the pullback from the Remmert reduction is strictly psh.
Pure dimension and continuity of the real-valued exhaustion are used
in the proof.  The broader convention allowing an upper
semicontinuous, possibly \(-\infty\)-valued exhaustion in
\cite{Vajaitu2022} is not treated here.
\end{remark}

In particular, a continuously weakly complete pure \(n\)-dimensional
space without a smooth psh exhaustion cannot admit a holomorphic
map to \(\C^{n-1}\) all of whose fibers have dimension at most one.
For the threefolds in Theorems~\ref{thm:main}
and~\ref{thm:nonkahler-threefold}, this is also explained directly
by their compact surfaces: the zero section in \(M\) and exceptional
surfaces in \(N\) are mapped to points by every holomorphic map to
\(\C^2\).

\begin{thebibliography}{99}
\raggedright

\bibitem{Koike2021}
T.~Koike,
\emph{Linearization of transition functions of a semi-positive line
bundle along a certain submanifold},
Ann. Inst. Fourier \textbf{71} (2021), no.~5, 2237--2271.
\href{https://doi.org/10.5802/aif.3439}{doi:10.5802/aif.3439}.

\bibitem{Wu2026}
X.~Wu,
\emph{Generalized Ueda Obstruction Classes and Applications to
Non-semi-positivity of Line Bundles},
preprint (2026), arXiv:2606.22821v1.
\href{https://arxiv.org/abs/2606.22821v1}{arXiv:2606.22821}.

\bibitem{Zhang2025}
Y.~Zhang,
\emph{On the Non-Semipositivity of a Nef and Big Line Bundle on
Grauert's Example},
preprint (2025), arXiv:2512.23446v1.
\href{https://arxiv.org/abs/2512.23446v1}{arXiv:2512.23446}.

\bibitem{Vajaitu2022}
V.~V\^{\i}j\^{\i}itu,
\emph{Holomorphic convexity and pseudoconvex surfaces},
Ann. Sc. Norm. Super. Pisa Cl. Sci. (5) \textbf{23} (2022),
no.~4, 1659--1678.
\href{https://doi.org/10.2422/2036-2145.202011_018}{doi:10.2422/2036-2145.202011\_018}.
\url{https://journals.sns.it/index.php/annaliscienze/article/view/5700}.

\bibitem{MongodiSlodkowski2020}
S.~Mongodi and Z.~Slodkowski,
\emph{Domains with a continuous exhaustion in weakly complete surfaces},
Math. Z. \textbf{296} (2020), nos.~3--4, 1011--1019.
\href{https://doi.org/10.1007/s00209-020-02466-z}{doi:10.1007/s00209-020-02466-z}.
\url{https://arxiv.org/abs/1904.03497}.

\bibitem{Grauert1962}
H.~Grauert,
\emph{\"Uber Modifikationen und exzeptionelle analytische Mengen},
Math. Ann. \textbf{146} (1962), 331--368.
\href{https://doi.org/10.1007/BF01441136}{doi:10.1007/BF01441136}.

\bibitem{BHPV2004}
W.~P.~Barth, K.~Hulek, C.~A.~M.~Peters and A.~Van de Ven,
\emph{Compact Complex Surfaces}, second enlarged edition,
Ergebnisse der Mathematik und ihrer Grenzgebiete (3), vol.~4,
Springer-Verlag, Berlin, 2004.
\href{https://doi.org/10.1007/978-3-642-57739-0}{doi:10.1007/978-3-642-57739-0}.

\bibitem{FN2005}
Y.~Fujimoto and N.~Nakayama,
\emph{Compact complex surfaces admitting non-trivial surjective endomorphisms},
Tohoku Math. J. (2) \textbf{57} (2005), no.~3, 395--426.
\href{https://doi.org/10.2748/tmj/1128703004}{doi:10.2748/tmj/1128703004}.
\href{https://www.kurims.kyoto-u.ac.jp/preprint/file/RIMS1431.pdf}{Author preprint, Section~2 and Lemma~2.1}.

\bibitem{Enoki}
I.~Enoki, \emph{Surfaces of class $\mathrm{VII}_0$ with curves}, Tohoku Math. J. (2)
\textbf{33} (1981), no.~4, 453--492.
\href{https://doi.org/10.2748/tmj/1178229349}{doi:10.2748/tmj/1178229349}.
\href{https://www.jstage.jst.go.jp/article/tmj1949/33/4/33_4_453/_pdf/-char/en}{Original article, pp.453--454}.

\bibitem{Fujino2012}
O.~Fujino,
\emph{Minimal model theory for log surfaces},
Publ. Res. Inst. Math. Sci. \textbf{48} (2012), no.~2, 339--371.
\href{https://doi.org/10.2977/PRIMS/71}{doi:10.2977/PRIMS/71}.

\bibitem{Fujino2021}
O.~Fujino,
\emph{Minimal model theory for log surfaces in Fujiki's class $\mathcal C$},
Nagoya Math. J. \textbf{244} (2021), 256--282.
\href{https://doi.org/10.1017/nmj.2020.14}{doi:10.1017/nmj.2020.14}.

\bibitem{Ueda1983}
T.~Ueda,
\emph{On the neighborhood of a compact complex curve with topologically
trivial normal bundle},
J. Math. Kyoto Univ. \textbf{22} (1982/83), no.~4, 583--607.
\href{https://doi.org/10.1215/kjm/1250521670}{doi:10.1215/kjm/1250521670}.

\bibitem{CLPT2018}
B.~Claudon, F.~Loray, J.~V.~Pereira and F.~Touzet,
\emph{Compact leaves of codimension one holomorphic foliations on
projective manifolds},
Ann. Sci. \'Ecole Norm. Sup. (4) \textbf{51} (2018), no.~6, 1457--1506.
\href{https://doi.org/10.24033/asens.2379}{doi:10.24033/asens.2379}.
\href{https://arxiv.org/abs/1512.06623}{Author preprint, Theorems~1.3, 2.5, 4.4 and Proposition~2.7}.

\bibitem{CFSTZ2026}
Y.~Chen, S.~Filip, S.~Sun, V.~Tosatti and J.~Zhang,
\emph{Compact K\"ahler surfaces with semipositive anticanonical bundle},
preprint (2026), arXiv:2609.26716v1, submitted September~22, 2026.
\href{https://arxiv.org/abs/2609.26716v1}{arXiv:2609.26716}.

\bibitem{KoikeUeda2025}
T.~Koike and T.~Ueda,
\emph{Plurisubharmonic functions on affine line bundles over compact
K\"ahler manifolds},
Publ. Res. Inst. Math. Sci. \textbf{61} (2025), no.~1, 139--152.
\href{https://doi.org/10.4171/PRIMS/61-1-3}{doi:10.4171/PRIMS/61-1-3}.

\bibitem{Barlet2015}
D.~Barlet,
\emph{Some examples due to H.~Hironaka},
preprint (2015), arXiv:1501.01129v2, Section~6.
\href{https://arxiv.org/abs/1501.01129}{arXiv:1501.01129}.

\bibitem{StacksHironaka}
The Stacks Project Authors,
\emph{The Stacks Project},
Lemma~110.67.1, Tag~08J1, ``Being projective is not local on the base.''
\url{https://stacks.math.columbia.edu/tag/08J1}.

\bibitem{DNS22}
T.-C.~Dinh, V.-A.~Nguy\^en, and N.~Sibony,
\emph{Unique ergodicity for foliations on compact K\"ahler surfaces},
Duke Math. J. \textbf{171} (2022), no.~13, 2627--2698.
\href{https://doi.org/10.1215/00127094-2022-0044}{doi:10.1215/00127094-2022-0044}.
Author version: \href{https://arxiv.org/abs/1811.07450v2}{arXiv:1811.07450v2}.
The numbered equations and Appendix~B cited here may be found in this
author version.

\bibitem{Wlodarczyk05}
J.~W{\l}odarczyk,
\emph{Simple Hironaka resolution in characteristic zero},
J. Amer. Math. Soc. \textbf{18} (2005), no.~4, 779--822.
\href{https://doi.org/10.1090/S0894-0347-05-00493-5}{doi:10.1090/S0894-0347-05-00493-5}.
Author version: \href{https://arxiv.org/abs/math/0401401}{arXiv:math/0401401}.


\bibitem{VijituForum2022}
V.~V\^{\i}j\^{\i}itu,
\emph{Holomorphic convexity of pseudoconvex spaces in terms of
the rank of structural sheaf},
Forum Math. \textbf{34} (2022), no.~2, 455--468.
\href{https://doi.org/10.1515/forum-2021-0262}{doi:10.1515/forum-2021-0262}.

\bibitem{BarletMagnusson1998}
D.~Barlet and J.~Magn\'usson,
\emph{Int\'egration de classes de cohomologie m\'eromorphes et
diviseurs d'incidence},
Ann. Sci. \'Ecole Norm. Sup. (4) \textbf{31} (1998),
no.~6, 811--842.
\href{https://doi.org/10.1016/S0012-9593(99)80004-4}
{doi:10.1016/S0012-9593(99)80004-4}.
\url{https://www.numdam.org/item/ASENS_1998_4_31_6_811_0.pdf}.

\bibitem{BarletMagnusson2014}
D.~Barlet and J.~Magn\'usson,
\emph{Cycles analytiques complexes I: th\'eor\`emes de
pr\'eparation des cycles},
Cours Sp\'ecialis\'es, vol.~22,
Soci\'et\'e Math\'ematique de France, Paris, 2014.

\bibitem{BarletMagnusson2019}
D.~Barlet and J.~Magn\'usson,
\emph{On nearly smooth complex spaces},
Math. Z. \textbf{292} (2019), 317--341.
\href{https://doi.org/10.1007/s00209-018-2202-2}{doi:10.1007/s00209-018-2202-2}.
\href{https://arxiv.org/abs/1710.08105}{Author preprint, arXiv:1710.08105}.

\bibitem{BarletAlpha2017}
D.~Barlet, \emph{The sheaf \(\alpha_X^\bullet\)},
preprint (2017), arXiv:1707.07962.
\href{https://arxiv.org/abs/1707.07962}{arXiv:1707.07962}.

\bibitem{Barlet1999}
D.~Barlet, \emph{How to use the cycle space in complex geometry},
in \emph{Several Complex Variables},
M.~Schneider and Y.-T.~Siu (eds.),
MSRI Publications, vol.~37, Cambridge University Press, 1999,
pp.~25--42.
\url{https://library.slmath.org/books/Book37/files/barlet.pdf}.

\end{thebibliography}
\end{document}
